% Optimization Methods and Software submission
% Taylor & Francis Interact layout (reference style S)
% Ported from the Elsevier cas-sc manuscript (optimization_ready.tex).
\documentclass[]{interact}

\usepackage{epstopdf}% To incorporate .eps illustrations using PDFLaTeX, etc.
\usepackage[caption=false]{subfig}% Support for small, `sub' figures and tables

\usepackage[numbers,sort&compress]{natbib}% Citation support using natbib.sty
\bibpunct[, ]{[}{]}{,}{n}{,}{,}% Citation support using natbib.sty
\renewcommand\bibfont{\fontsize{10}{12}\selectfont}% Bibliography support using natbib.sty

% --- Packages carried over from the source manuscript ---
\usepackage{amssymb}
\usepackage{float}
\usepackage{algorithm}
\usepackage{multirow}
\usepackage{array}
\usepackage{longtable}
\usepackage{enumitem}
\usepackage{xcolor,url}\graphicspath{{figures/}}
\restylefloat{figure}

% --- Theorem-like environments ---
\theoremstyle{plain}
\newtheorem{theorem}{Theorem}[section]
\newtheorem{lemma}[theorem]{Lemma}
\newtheorem{corollary}[theorem]{Corollary}
\newtheorem{proposition}[theorem]{Proposition}

\theoremstyle{definition}
\newtheorem{definition}[theorem]{Definition}
\newtheorem{example}[theorem]{Example}
\newtheorem{assumption}{Assumption}

\theoremstyle{remark}
\newtheorem{remark}{Remark}
\newtheorem{notation}{Notation}

% --- Author macros carried over from the source manuscript ---
\newcommand{\algstep}[1]{\smallskip\noindent\textbf{#1}}
\newcommand{\R}{\mathbb{R}}
\newcommand{\N}{\mathbb{N}}
\newcommand{\norm}[1]{\left\lVert #1 \right\rVert}
\newcommand{\abs}[1]{\left\lvert #1 \right\rvert}
\newcommand{\inner}[2]{\left\langle #1,\, #2 \right\rangle}
\DeclareMathOperator*{\argmin}{arg\,min}
\DeclareMathOperator*{\argmax}{arg\,max}
\DeclareMathOperator{\conv}{conv}
\DeclareMathOperator{\diag}{diag}
\DeclareMathOperator{\vecop}{vec}
\newcommand{\cvx}{\mathrm{cvx}}
\newcommand{\ps}{\mathrm{ps}}
\newcommand{\dtap}{\textsc{DTAP}}

\begin{document}

\title{A Dual-Track Adaptive Projection Algorithm for a Nonlinear Split Feasibility Problem with Mixed Convex--Pseudoconvex Constraints}

\author{
\name{Nguyen Van The\textsuperscript{a}\thanks{Email: the.nv227153@sis.hust.edu.vn}, Tran Ngoc Thang\textsuperscript{a}\thanks{Email: thang.tranngoc@hust.edu.vn} and Nguyen Hong Son\textsuperscript{a}\thanks{CONTACT Nguyen Hong Son. Email: Son.NH240066D@sis.hust.edu.vn}}
\affil{\textsuperscript{a}Faculty of Mathematics and Informatics, Hanoi University of Science and Technology, No.~1 Dai Co Viet Street, Bach Mai Ward, Hanoi 100000, Vietnam}
}

\maketitle

\begin{abstract}
We study a nonlinear split feasibility problem \leavevmode\textcolor{red}{in $\R^n$} whose output set is a normal box and whose components \leavevmode\textcolor{red}{are differentiable and either convex or} pseudoconvex. \leavevmode\textcolor{red}{The Dual-Track Adaptive Projection (DTAP) algorithm splits the constraints into a convex and a pseudoconvex track. It takes a residual-scaled Polyak-type step on the nearly most violated convex constraints and a safeguarded diminishing step on the nearly most violated pseudoconvex ones, combines the two steps by a proportional weight, and projects onto the base set. The weight is a residual-responsive choice, without a claim of optimality. The step sizes require no prescribed operator norm, Lipschitz constant, or H\"older modulus; the proof uses bounds over a bounded set containing the iterates. If the problem is consistent and the components are continuously differentiable, the whole sequence converges to a feasible point. In a feature-selection application all three constraints are convex. Experiments with constraint-activity counts, tuned subgradient methods and nonconvex pseudoconvex test problems show that tuned subgradient methods can match or outperform DTAP, and that DTAP is sensitive to its starting point and step scale.}
\end{abstract}

\begin{keywords}
Split feasibility problem; Pseudoconvex optimization; Tchebyshev scalarization; Adaptive step size; \leavevmode\textcolor{red}{Global} convergence; Feature selection
\end{keywords}
\begin{amscode}
47J25; 90C25; 90C26; 65K05; 68T05
\end{amscode}

\section{Introduction}\label{sec:intro}

\leavevmode\textcolor{red}{The split feasibility problem (SFP) was introduced by Censor and Elfving~\citep{censor1994} as a model for constrained inverse problems. Given nonempty closed convex sets $C\subseteq\R^n$ and $Q\subseteq\R^m$ and a matrix $\mathcal A\in\R^{m\times n}$, it asks to}
\begin{equation}\label{eq:sfp-intro}
\text{find }x^\star\in C\quad\text{such that}\quad\mathcal A x^\star\in Q.
\end{equation}
\leavevmode\textcolor{red}{Byrne's $CQ$ algorithm~\citep{byrne2002} iterates}
\begin{equation}\label{eq:cq-intro}\textcolor{red}{
x^{k+1}=P_C\bigl(x^k-\tau\mathcal A^\top(I-P_Q)\mathcal A x^k\bigr),
\qquad 0<\tau<2/\norm{\mathcal A}^2.
}\end{equation}
\leavevmode\textcolor{red}{Here $\norm{\mathcal A}$ is the spectral norm of $\mathcal A$. Relaxed variants project onto supporting half-spaces~\citep{yang2004}; self-adaptive steps avoid estimating $\norm{\mathcal A}$~\citep{zhang2009adaptive,zhao2011adaptive,lopez2012}; operator splitting~\citep{zhang2012simultaneous} and alternating direction methods~\citep{yang2014alternating} decompose the problem. These ideas are standard; DTAP adapts them to the nonlinear case.}

\leavevmode\textcolor{red}{Replacing $\mathcal A$ by a nonlinear map $F$ gives the nonlinear SFP~\citep{li2013nsfp,yen2019}}
\begin{equation}\label{eq:nsfp-intro}
\text{find }x^\star\in K\quad\text{such that}\quad F(x^\star)\in Q,
\end{equation}
\leavevmode\textcolor{red}{where $K\subseteq\R^n$ is nonempty, closed and convex, $Q=\prod_{i=1}^m(-\infty,b_i]$ is a normal box, and each $F_i$ is differentiable and convex or pseudoconvex on $K$. Since pseudoconvex functions have convex sublevel sets, the feasible set $\Omega=\{x\in K:F(x)\in Q\}$ is convex. We seek a point of $\Omega$, not a minimizer of any $F_i$}.

\leavevmode\textcolor{red}{Li et al.~\citep{li2013nsfp} assume a convex aggregate proximity function, which componentwise pseudoconvexity does not ensure. Yen et al.~\citep{yen2019} use the same max-type violation measure and active-set device, but require quasilinear components, a class incomparable with ours. Hu et al.~\citep{hu2022quasi} treat quasiconvex inequalities, which contain our class, without using the gradient inequality of convex components.}

\leavevmode\textcolor{red}{Because $Q$ is a box, the $\ell_\infty$-distance from $F(x)$ to $Q$ has the closed form}
\[
J(x)=d_\infty(F(x),Q)=\max_{1\le i\le m}(F_i(x)-b_i)^+.
\]
\leavevmode\textcolor{red}{DTAP splits the constraints into a convex track and a pseudoconvex track. On each track it sums the gradients of the nearly most violated constraints, weighted by their violations. The gradient inequality of convex functions justifies a residual-scaled Polyak-type step on the convex track~\citep{polyak1969,lopez2012}; pseudoconvexity gives only the sign of a directional derivative, so the pseudoconvex track uses a safeguarded diminishing step. A proportional weight combines the two steps, and the result is projected onto $K$; the proof uses this weight only to keep a persistently violated track active (Remark~\ref{rem:weight-scope}). No prescribed operator norm, Lipschitz constant, or H\"older modulus is needed: a quasi-Fej\'er inequality shows that the iterates are bounded, and the proof then uses local bounds, without Lipschitz gradients.}

\leavevmode\textcolor{red}{As an application, upper bounds on three criteria from a fractional model~\citep{wang2021fractional} and a fuzzy-rough model~\citep{wang2026afrsfcv} define the constraints; a feasible weight vector satisfies all three bounds but need not solve either source problem. These constraints are convex (Proposition~\ref{prop:fs-reg}), so we also use test problems with a nonconvex pseudoconvex constraint to exercise the pseudoconvex track.}

\leavevmode\textcolor{red}{The contributions are as follows.}
\begin{enumerate}[label=\textbf{C\arabic*.},leftmargin=2.4em,itemsep=3pt]
\item \leavevmode\textcolor{red}{A dual-track method that uses the gradient inequality on the convex constraints and only sign information on the pseudoconvex ones, without assuming a convex proximity function, unlike~\citep{li2013nsfp}, or a H\"older modulus, unlike the dynamic step of~\citep{hu2022quasi}}.
\item \leavevmode\textcolor{red}{Convergence of the whole sequence to a feasible point under Assumption~\ref{ass:standing} (Theorem~\ref{thm:strong}}).
\item \leavevmode\textcolor{red}{Experiments that report constraint activity, compare DTAP with tuned subgradient methods, and measure its sensitivity to the starting point and the step scale.}
\end{enumerate}

Section~\leavevmode\textcolor{red}{\ref{sec:prelim} gives preliminaries. Section~\ref{sec:convergence} presents} \dtap{} and its convergence analysis. Section~\leavevmode\textcolor{red}{\ref{sec:experiments} reports experiments, and Section~\ref{sec:conclusion} concludes.}


\section{Preliminaries}\label{sec:prelim}
Throughout, $H=\R^n$ \leavevmode\textcolor{red}{is} the Euclidean space with inner product \leavevmode\textcolor{red}{$\inner{\cdot}{\cdot}$ and induced norm $\norm{\cdot}$. All convergence statements concern the norm topology. Bounded sequences have convergent subsequences. The standard projection and generalized convexity facts below can be found in~\citep{bauschke2011,mangasarian1994}}.

\leavevmode\textcolor{red}{For a nonempty closed convex set $K\subseteq H$, the metric projection $P_K(x)$ is the unique minimizer of $\norm{y-x}$ over $y\in K$}.
\begin{proposition}\label{prop:proj-char}
For $x\in H$ and $p\in K$, $p=P_K(x)$ if and only if
\begin{equation}\label{eq:proj-char}\textcolor{red}{
\inner{x-p}{y-p}\le0\qquad\text{for every }y\in K.
}\end{equation}
\end{proposition}
\leavevmode\textcolor{red}{This characterization yields Lemma~\ref{lem:proj-nonexp} below.}

\leavevmode\textcolor{red}{A function $f:K\to\R$ is convex if
$f((1-t)x+ty)\le(1-t)f(x)+tf(y)$ for $x,y\in K$ and $t\in[0,1]$.
The following weaker derivative condition was introduced by Mangasarian~\citep{mangasarian1965}}.
\begin{definition}\label{def:pseudoconvex}
\leavevmode\textcolor{red}{A function $f$ that is differentiable on an open neighbourhood of $K$ is pseudoconvex on $K$ if, for all $x,y\in K$},
\begin{equation}\label{eq:contrapositive}
f(y)<f(x)\quad\Longrightarrow\quad\inner{\nabla f(x)}{y-x}<0.
\end{equation}
\end{definition}
\leavevmode\textcolor{red}{Every differentiable convex function is pseudoconvex. Pseudoconvex functions have convex sublevel sets, although the functions themselves need not be convex.}

\section{The DTAP algorithm and its convergence analysis}\label{sec:convergence}
\leavevmode\textcolor{red}{We formulate the normal-box NSFP, state DTAP, and prove convergence of the whole sequence
to a feasible point.}

\subsection{Problem statement and Tchebyshev scalarization}\label{ssec:model-stmt}
\leavevmode\textcolor{red}{Let $K\subseteq H$ be nonempty closed convex, and let $U\subseteq H$ be an open
neighbourhood of $K$}. Fix an integer $m\ge1$, let \leavevmode\textcolor{red}{$F=(F_1,\dots,F_m):U\to\R^m$},
let $b=(b_1,\dots,b_m)\in\R^m$, and set the normal box
$Q:=(-\infty,b_1]\times\cdots\times(-\infty,b_m]$. We study the instance
of~\eqref{eq:nsfp-intro}
\begin{equation}\label{eq:nsfp-box}
\text{find } x^\star\in\Omega,\qquad
\Omega:=\bigl\{x\in K: F_i(x)\le b_i,\ i=1,\dots,m\bigr\},
\end{equation}
and call $\Omega$ the \emph{feasible set}. To turn $F(x)\in Q$ into a tractable merit
function, we measure the residual in the sup-norm. For $z=F(x)$,
$d_\infty(z,Q)=\inf_{q\le b}\max_{1\le i\le m}\abs{z_i-q_i}$. By the box structure the
infimum separates over coordinates: if $z_i\le b_i$ then
$\inf_{q_i\le b_i}\abs{z_i-q_i}=0$ (at $q_i=z_i$), and if $z_i>b_i$ then it equals
$z_i-b_i$ (at $q_i=b_i$). Hence $\inf_{q_i\le b_i}\abs{z_i-q_i}=(z_i-b_i)^+$ with
$(t)^+:=\max\{t,0\}$, so
\begin{equation}\label{eq:dinfty}
d_\infty(z,Q)=\max_{1\le i\le m}(z_i-b_i)^+ .
\end{equation}
Define the \emph{error function}
\begin{equation}\label{eq:Jdef}
J(x):=d_\infty\bigl(F(x),Q\bigr)=\max_{1\le i\le m}\bigl(F_i(x)-b_i\bigr)^+ .
\end{equation}
\leavevmode\textcolor{red}{Then $x\in K$ is feasible if and only if $J(x)=0$; when $\Omega\ne\emptyset$, the feasible
points are the minimizers of $J$ over $K$. Formula~\eqref{eq:Jdef} gives a closed-form feasibility residual whose
maximizers identify the most violated constraints. Section~\ref{ssec:dual-track}
introduces an approximate active set to handle the nonsmooth maximum.}

\subsection{\textcolor{red}{Decomposition into two tracks}}\label{ssec:dual-track}
\leavevmode\textcolor{red}{We split $I:=\{1,\dots,m\}$ into a \emph{convex track} $I^{\cvx}$ and a
\emph{pseudoconvex track} $I^{\ps}$, each with its own step rule. The convex track uses a Polyak step~\citep{polyak1969};
the pseudoconvex track uses a safeguarded diminishing step.
A proportional weight combines the two updates.}

Fix a parameter $\delta>0$. For \leavevmode\textcolor{red}{$x\in U$} define the componentwise violations
\begin{equation}\label{eq:vi}\textcolor{red}{
v_i(x):=\bigl(F_i(x)-b_i\bigr)^+,
}\end{equation}
\leavevmode\textcolor{red}{so that $J(x)=\max_i v_i(x)$ by~\eqref{eq:Jdef}. Assumption~\ref{ass:standing} makes the
partition precise.} The \emph{per-track maximal violations} are
\begin{align}
J^{\cvx}(x)&:=\max_{i\in I^{\cvx}}v_i(x)\quad(\,{=}\,0\text{ if }I^{\cvx}=\emptyset),\label{eq:Jcvx}\\
J^{\ps}(x)&:=\max_{j\in I^{\ps}}v_j(x)\quad(\,{=}\,0\text{ if }I^{\ps}=\emptyset),\label{eq:Jps}
\end{align}
so $J(x)=\max\{J^{\cvx}(x),J^{\ps}(x)\}$. To cope with the nonsmoothness of the maximum,
we replace the \leavevmode\textcolor{red}{set of indices at which the maximum is attained} by a $\delta$-enlarged
\emph{approximate active set} on each track \leavevmode\textcolor{red}{(`active' meaning within $\delta$ of the largest
violation):}
\begin{align}
K^{\cvx}(x)&:=\begin{cases}\{i\in I^{\cvx}:v_i(x)\ge J^{\cvx}(x)-\delta\} & \text{if }J^{\cvx}(x)>0,\\ \emptyset & \text{if }J^{\cvx}(x)=0,\end{cases}\label{eq:Kcvx}\\[2pt]
K^{\ps}(x)&:=\begin{cases}\{j\in I^{\ps}:v_j(x)\ge J^{\ps}(x)-\delta\} & \text{if }J^{\ps}(x)>0,\\ \emptyset & \text{if }J^{\ps}(x)=0.\end{cases}\label{eq:Kps}
\end{align}
On these sets we form the \emph{per-track local functions and directions}
\begin{align}\color{red}
\Phi^{\cvx}(x)&\color{red}:=\tfrac12\!\!\sum_{i\in K^{\cvx}(x)}\!\! v_i(x)^2, &\color{red} g^{\cvx}(x)&\color{red}:=\!\!\sum_{i\in K^{\cvx}(x)}\!\! v_i(x)\nabla F_i(x),\label{eq:phi-g-cvx}\\\color{red}
&\color{red} &\color{red} g^{\ps}(x)&\color{red}:=\!\!\sum_{j\in K^{\ps}(x)}\!\! v_j(x)\nabla F_j(x).\label{eq:phi-g-ps}
\end{align}
The vectors $g^{\cvx}(x),g^{\ps}(x)$ are the two search directions, and $\Phi^{\cvx}(x)$
supplies the Polyak numerator on the convex track. Along the iterates we abbreviate
$J^{\cvx}_k:=J^{\cvx}(x^k)$, $g^{\cvx}_k:=g^{\cvx}(x^k)$, $\Phi^{\cvx}_k:=\Phi^{\cvx}(x^k)$,
and likewise on the pseudoconvex track.

\subsection{\textcolor{red}{Assumptions}}\label{ssec:assumptions}
\begin{assumption}\label{ass:standing}
We assume \textup{(A1)--(A3)}:
\begin{itemize}[leftmargin=2.4em,itemsep=2pt]
\item[\textup{(A1)}] \textup{(a)} $\Omega\ne\emptyset$ \leavevmode\textcolor{red}{(consistency).} \textup{(b)} $I=\{1,\dots,m\}$
is partitioned as $I=I^{\cvx}\cup I^{\ps}$; each \leavevmode\textcolor{red}{$F_i:U\to\R$ is convex on $K$} if
$i\in I^{\cvx}$ and pseudoconvex \leavevmode\textcolor{red}{on $K$} if $i\in I^{\ps}$.
\item[\textup{(A2)}] \leavevmode\textcolor{red}{Each $F_i:U\to\R$ is continuously differentiable on $U$}.
\item[\textup{(A3)}] \textup{(i)} $\delta>0$ fixed. \textup{(ii)}
$\rho_{\cvx}\in(0,2)$. \textup{(iii)} $\{\gamma_k\}\subset(0,\infty)$ with
$\sum_{k=0}^\infty\gamma_k^2<\infty$ and $\sum_{k=0}^\infty\gamma_k=\infty$.
\end{itemize}
\end{assumption}
\leavevmode\textcolor{red}{The open-neighbourhood formulation permits fractional functions whose denominators are
positive on $K$; the algorithm evaluates the functions and gradients only at points of $K$.
Although some component functions may be nonconvex, $\Omega$ is convex: differentiable
pseudoconvex functions are quasiconvex on $K$, so all constraint sublevel sets in $K$ are convex.}
\begin{example}\label{ex:gamma}
The sequence $\gamma_k:=1/(k+1)$ satisfies (A3)(iii), since $\sum_{k\ge0}(k+1)^{-2}<\infty$
while $\sum_{k\ge0}(k+1)^{-1}=\infty$. More generally $\gamma_k=(k+c)^{-s}$ with $c>0$
and $s\in(1/2,1]$ qualifies; the experiments use $s=0.55$.
\end{example}

\subsection{The DTAP algorithm}\label{ssec:algorithm}
The Dual-Track Adaptive Projection algorithm combines the two track-specific step sizes
through the \leavevmode\textcolor{red}{proportional} weight $\mu_k$ and projects the combined update onto the base convex
set $K$. It is stated as Algorithm~\ref{alg:dtap}.

\begin{algorithm}[H]
\caption{Dual-Track Adaptive Projection (DTAP) for the normal-box NSFP
\eqref{eq:nsfp-box}.}
\label{alg:dtap}
\small
\noindent\textbf{Input:} map $F=(F_1,\dots,F_m)$, thresholds $b$, base set $K$, index
split $I^{\cvx}\cup I^{\ps}$.\\
\textbf{Parameters:} $\delta>0$, tolerance $\varepsilon>0$, $\rho_{\cvx}\in(0,2)$,
sequence $\{\gamma_k\}$ satisfying (A3)(iii).\\
\textbf{Initialization:} choose $x^0\in K$; set $k\leftarrow 0$.

\algstep{Step 1 (per-track quantities).} If $J(x^k)\le\varepsilon$, stop and return $x^k$; otherwise:
\begin{itemize}[leftmargin=1.6em,itemsep=1pt,topsep=2pt]
\item[(1a)] Compute $v_i(x^k)$ for all $i$, and $J^{\cvx}_k=\max_{i\in I^{\cvx}}v_i(x^k)$,
$J^{\ps}_k=\max_{j\in I^{\ps}}v_j(x^k)$ (each $0$ if its track is empty).
\item[(1b)] Form the approximate active sets $K^{\cvx}_k,K^{\ps}_k$ via
\eqref{eq:Kcvx}--\eqref{eq:Kps}.
\item[(1c)] Form $\Phi^{\cvx}_k$, $g^{\cvx}_k$, $g^{\ps}_k$ via
\eqref{eq:phi-g-cvx}--\eqref{eq:phi-g-ps}.
\item[(1d)] Per-track step sizes (Polyak on the convex track, safeguarded on the
pseudoconvex track):
\[
\eta^{\cvx}_k=\begin{cases}\rho_{\cvx}\,\dfrac{2\Phi^{\cvx}_k}{\norm{g^{\cvx}_k}^2} & \text{if }\norm{g^{\cvx}_k}>0,\\ 0 & \text{otherwise,}\end{cases}
\qquad
\eta^{\ps}_k=\frac{\gamma_k}{\max\{1,\norm{g^{\ps}_k}\}} .
\]
\item[(1e)] Proportional weight:
$\displaystyle \mu_k=\frac{J^{\cvx}_k}{J^{\cvx}_k+J^{\ps}_k}$ if
$J^{\cvx}_k+J^{\ps}_k>0$, and $\mu_k=1/2$ otherwise.
\item[(1f)] Combine the two directions:
$\displaystyle y^k=x^k-\bigl[\mu_k\eta^{\cvx}_k g^{\cvx}_k+(1-\mu_k)\eta^{\ps}_k g^{\ps}_k\bigr].$
\end{itemize}
\algstep{Step 2 (projection).}\quad \leavevmode\textcolor{red}{$x^{k+1}=P_K(y^k)$}; set $k\leftarrow k+1$ and go to Step~1.
\end{algorithm}

\begin{remark}\label{rem:weight-scope}
\leavevmode\textcolor{red}{The proportional weight is an adaptive allocation based on the two maximal violations;
we do not claim that it is the only choice for which DTAP converges, or that it is optimal.
Its specific form is used in~\eqref{eq:p10} to control the convex-track residual and in
Claim~2 of Stage~2 of Theorem~\ref{thm:strong} to ensure that a persistently violated
pseudoconvex track retains positive weight. Although the one-step inequality of Lemma~\ref{lem:fejer} holds
for every $\mu_k\in[0,1]$, it does not by itself imply convergence for every weight sequence.}
\end{remark}

\begin{remark}\label{rem:regularity-scope}
\leavevmode\textcolor{red}{The iteration does not require prescribed operator norms, Lipschitz constants, or
H\"older moduli. The proof nevertheless uses finite bounds on function values and
gradients on a compact set containing the iterates, after boundedness has been proved.
Bounded gradients imply Lipschitz continuity of the functions on that convex compact set,
but these bounds are not algorithm inputs. Assumption (A2) does not
imply Lipschitz continuity of the gradients, and the theorem does not require it.
The user still chooses $\rho_{\cvx}$, $\delta$, and $\{\gamma_k\}$; these parameters do not
depend on any constant of the problem.}
\end{remark}

For the analysis we write the combined increment as
\begin{equation}\label{eq:dk}
d_k:=\mu_k\eta^{\cvx}_k g^{\cvx}_k+(1-\mu_k)\eta^{\ps}_k g^{\ps}_k,
\end{equation}
so $y^k=x^k-d_k$ and $x^{k+1}=P_K(x^k-d_k)$.
\subsection{Per-iteration cost}\label{ssec:complexity}
Each \dtap{} iteration evaluates the
componentwise violations $v_i(x^k)$ and the active-set gradients $\{\nabla F_i(x^k):i\in
K^{\cvx}_k\cup K^{\ps}_k\}$, forms the two search directions $g^{\cvx}_k,g^{\ps}_k$ as weighted
sums of these gradients, computes the step sizes, and \leavevmode\textcolor{red}{performs one projection onto $K$.
It uses neither an operator-norm estimate nor a line search. The cost is dominated by
the constraint evaluations and the projection. The cost for the feature-selection instance
is given in Section~\ref{sec:fs}
(Remark~\ref{rem:fs-cost}}).

\subsection{Auxiliary lemmas}\label{ssec:aux}
\leavevmode\textcolor{red}{The two tracks need different arguments. On the convex track, Lemma~\ref{lem:grad-cvx} bounds
$\inner{\nabla F_i(x)}{x-x^\star}$ from below by the violation $v_i(x)$; this gives a decrease of
$\norm{x^k-x^\star}^2$ of order $\mu_kT_k$. On the pseudoconvex track,
definition~\eqref{eq:contrapositive} gives only the sign of this inner product.
Lemma~\ref{lem:fejer} combines both tracks in one quasi-Fej\'er inequality; the choice of the
parameter $\zeta$ in Young's inequality keeps the convex-track term negative. Stage~2 of
Theorem~\ref{thm:strong} then shows, by contradiction, that at least one cluster point is
feasible, and Stage~3 shows that the whole sequence converges to it. The two degenerate cases
follow as corollaries.}

Throughout, $\{x^k\}$ is the infinite sequence generated by Algorithm~\ref{alg:dtap} with
the stopping test removed. As shown below
(Lemma~\ref{lem:jcvx0} and Theorem~\ref{thm:strong}), $J(x^k)\to 0$; hence for any
$\varepsilon>0$ there is a finite $k$ with $J(x^k)\le\varepsilon$, so the finite-stopping
version terminates after finitely many steps with an $\varepsilon$-feasible point.

We begin with two preparatory results used repeatedly below.

\begin{lemma}\label{lem:proj-nonexp}
Let $K\subseteq H$ be nonempty closed convex. For every $y\in H$ and $z\in K$,
\begin{equation}\label{eq:proj-nonexp}
\norm{P_K(y)-z}^2\le\norm{y-z}^2-\norm{y-P_K(y)}^2\le\norm{y-z}^2 .
\end{equation}
\end{lemma}
\begin{proof}
Put $p:=P_K(y)$. By Proposition~\ref{prop:proj-char}, $\inner{y-p}{z-p}\le0$, so expanding
$\norm{y-z}^2=\norm{y-p}^2+2\inner{y-p}{p-z}+\norm{p-z}^2$ and using $\inner{y-p}{p-z}\ge0$
gives $\norm{y-z}^2\ge\norm{y-p}^2+\norm{p-z}^2$, which is~\eqref{eq:proj-nonexp}.
\end{proof}

\begin{lemma}\label{lem:seq}
Let $\{a_k\}\subset[0,\infty)$ and $\{\varepsilon_k\}\subset[0,\infty)$ satisfy
$a_{k+1}\le a_k+\varepsilon_k$ for all $k$ and $\sum_{k=0}^\infty\varepsilon_k<\infty$.
Then $\{a_k\}$ converges to a finite limit.
\end{lemma}
\begin{proof}
$b_k:=a_k-\sum_{l<k}\varepsilon_l$ is nonincreasing (from $a_{k+1}\le a_k+\varepsilon_k$) and
bounded below by $-\sum_l\varepsilon_l>-\infty$, hence converges; since
$\sum_{l<k}\varepsilon_l\to S<\infty$, $a_k=b_k+\sum_{l<k}\varepsilon_l$ converges.
\end{proof}

On the convex track, convexity gives a quantitative link between the gradient and the
violation.

\begin{lemma}\label{lem:grad-cvx}
For $x^\star\in\Omega$, $i\in I^{\cvx}$, and $x\in K$,
\[
v_i(x)\,\inner{\nabla F_i(x)}{x-x^\star}\ge v_i(x)^2 .
\]
\end{lemma}
\begin{proof}
If $v_i(x)=0$ both sides vanish. Let $v_i(x)>0$, i.e.\ $F_i(x)>b_i$. Define
$\varphi:[0,1]\to\R$ by $\varphi(t):=F_i((1-t)x+tx^\star)$; since $K$ is convex,
$[x,x^\star]\subset K$, so $\varphi$ is well defined, and by (A1)(b) it is convex, while
by (A2) it is differentiable with $\varphi'(0^+)=\inner{\nabla F_i(x)}{x^\star-x}$. The
one-dimensional gradient inequality gives $\varphi(1)\ge\varphi(0)+\varphi'(0^+)$, i.e.
$F_i(x^\star)\ge F_i(x)+\inner{\nabla F_i(x)}{x^\star-x}$. As $x^\star\in\Omega$ we have
$F_i(x^\star)\le b_i<F_i(x)$, whence
$\inner{\nabla F_i(x)}{x-x^\star}\ge F_i(x)-b_i=v_i(x)$. Multiplying by $v_i(x)>0$ gives
the claim.
\end{proof}

\begin{remark}\label{rem:grad-cvx-sum}
Summing Lemma~\ref{lem:grad-cvx} over the approximate active set $K^{\cvx}_k$ yields
$\inner{g^{\cvx}_k}{x^k-x^\star}\ge\sum_{i\in K^{\cvx}_k}v_i(x^k)^2=2\Phi^{\cvx}_k$, a bound
used in the proofs of Lemmas~\ref{lem:fejer} and~\ref{lem:jcvx0}.
\end{remark}

On the pseudoconvex track, only a sign is available.

\begin{lemma}\label{lem:sign-ps}
For $x^\star\in\Omega$, $j\in I^{\ps}$, and $x\in K$,
$v_j(x)\,\inner{\nabla F_j(x)}{x-x^\star}\ge 0$, with strict inequality when $v_j(x)>0$.
\end{lemma}
\begin{proof}
If $v_j(x)=0$ both sides vanish. If $v_j(x)>0$ then $F_j(x)>b_j\ge F_j(x^\star)$. By
(A1)(b), $F_j|_K$ is pseudoconvex; applying~\eqref{eq:contrapositive}
with $y=x^\star$ gives $\inner{\nabla F_j(x)}{x^\star-x}<0$, i.e.
$\inner{\nabla F_j(x)}{x-x^\star}>0$. Multiplying by $v_j(x)>0$ yields the strict
inequality.
\end{proof}

Finally, two compactness-based norm bounds on the search directions.

\begin{lemma}\label{lem:bound-cvx}
Assume $I^{\cvx}\ne\emptyset$ and let \leavevmode\textcolor{red}{$B\subset U$} be compact. With
$L_B:=\max_{i\in I^{\cvx}}\sup_{x\in B}\norm{\nabla F_i(x)}$ (finite by (A2)), for every
$x\in B$,
\[
\norm{g^{\cvx}(x)}^2\le 2\,\abs{I^{\cvx}}\,L_B^2\,\Phi^{\cvx}(x).
\]
\end{lemma}
\begin{proof}
Put $K'=K^{\cvx}(x)$, $m'=\abs{K'}\le\abs{I^{\cvx}}$. By the triangle and Cauchy--Schwarz
inequalities,
\[
\norm{g^{\cvx}(x)}=\Bigl\|\sum_{i\in K'}v_i(x)\nabla F_i(x)\Bigr\|
\le\Bigl(\sum_{i\in K'}v_i(x)^2\Bigr)^{1/2}\Bigl(\sum_{i\in K'}\norm{\nabla F_i(x)}^2\Bigr)^{1/2}.
\]
Since $\sum_{i\in K'}v_i(x)^2=2\Phi^{\cvx}(x)$ and
$\sum_{i\in K'}\norm{\nabla F_i(x)}^2\le m'L_B^2\le\abs{I^{\cvx}}L_B^2$, squaring gives the
claim.
\end{proof}

\begin{lemma}\label{lem:bound-ps}
Assume $I^{\ps}\ne\emptyset$ and let \leavevmode\textcolor{red}{$B\subset U$} be compact. With
$L'_B:=\max_{j\in I^{\ps}}\sup_{x\in B}\norm{\nabla F_j(x)}$ and
$V_B:=\max_{j\in I^{\ps}}\sup_{x\in B}v_j(x)$ (both finite by (A2)), for every $x\in B$,
\[
\norm{g^{\ps}(x)}\le\abs{I^{\ps}}\,V_B\,L'_B .
\]
\end{lemma}
\begin{proof}
By the triangle inequality,
$\norm{g^{\ps}(x)}=\bigl\|\sum_{j\in K^{\ps}(x)}v_j(x)\nabla F_j(x)\bigr\|
\le\sum_{j\in K^{\ps}(x)}v_j(x)\norm{\nabla F_j(x)}\le\abs{K^{\ps}(x)}V_BL'_B
\le\abs{I^{\ps}}V_BL'_B$.
\end{proof}

\subsection{Boundedness and decay of the iteration error}\label{ssec:bounded}
\leavevmode\textcolor{red}{The following inequality is of quasi-Fej\'er type in the sense of
Combettes~\citep{combettes2001quasifejer}. It yields boundedness of $\{x^k\}$ and decay of the
convex-track violation.} Set
\begin{equation}\label{eq:c0C0}
c_0:=\frac{\rho_{\cvx}(2-\rho_{\cvx})}{2}\in\Bigl(0,\tfrac12\Bigr],\qquad
C_0:=\frac{2+\rho_{\cvx}}{2-\rho_{\cvx}}\in(1,\infty),
\end{equation}
and
\begin{equation}\label{eq:Tk}
T_k:=\begin{cases}\dfrac{(2\Phi^{\cvx}_k)^2}{\norm{g^{\cvx}_k}^2} & \text{if }\norm{g^{\cvx}_k}>0,\\ 0 & \text{if }\norm{g^{\cvx}_k}=0.\end{cases}
\end{equation}
For fixed $x^\star\in\Omega$ write $R^{x^\star}_k:=\inner{g^{\ps}_k}{x^k-x^\star}$.

\begin{lemma}\label{lem:fejer}
For every $x^\star\in\Omega$ and $k\ge 0$,
\begin{equation}\label{eq:fejer}
\norm{x^{k+1}-x^\star}^2\le\norm{x^k-x^\star}^2-c_0\mu_k T_k-2(1-\mu_k)\eta^{\ps}_k R^{x^\star}_k+C_0\gamma_k^2 .
\end{equation}
\end{lemma}
\begin{proof}
Apply Lemma~\ref{lem:proj-nonexp} with $y=y^k$ and $z=x^\star\in K$ and use
$y^k=x^k-d_k$ from~\eqref{eq:dk}:
\begin{equation}\label{eq:p1}
\norm{x^{k+1}-x^\star}^2\le\norm{x^k-x^\star}^2-2\inner{d_k}{x^k-x^\star}+\norm{d_k}^2 .
\end{equation}
\emph{Cross term.} Expanding $d_k$,
\begin{equation}\label{eq:p2}
\inner{d_k}{x^k-x^\star}=\mu_k\eta^{\cvx}_k\inner{g^{\cvx}_k}{x^k-x^\star}+(1-\mu_k)\eta^{\ps}_k R^{x^\star}_k .
\end{equation}
If $\norm{g^{\cvx}_k}=0$ then $\eta^{\cvx}_k=0$ and the first term vanishes. If
$\norm{g^{\cvx}_k}>0$, Remark~\ref{rem:grad-cvx-sum} gives
$\inner{g^{\cvx}_k}{x^k-x^\star}\ge 2\Phi^{\cvx}_k$; multiplying by
$\mu_k\eta^{\cvx}_k\ge 0$ and substituting
$\eta^{\cvx}_k=2\rho_{\cvx}\Phi^{\cvx}_k/\norm{g^{\cvx}_k}^2$,
\begin{equation}\label{eq:p4}
\mu_k\eta^{\cvx}_k\inner{g^{\cvx}_k}{x^k-x^\star}\ge\mu_k\rho_{\cvx}\frac{(2\Phi^{\cvx}_k)^2}{\norm{g^{\cvx}_k}^2}=\mu_k\rho_{\cvx}T_k .
\end{equation}
\emph{Quadratic term.} By Young's inequality
\leavevmode\textcolor{red}{$\norm{u+v}^2\le(1+\zeta)\norm{u}^2+(1+\zeta^{-1})\norm{v}^2$} with
$u=\mu_k\eta^{\cvx}_k g^{\cvx}_k$, $v=(1-\mu_k)\eta^{\ps}_k g^{\ps}_k$,
\[\textcolor{red}{
\norm{d_k}^2\le(1+\zeta)\rho_{\cvx}^2\mu_k^2 T_k+(1+\zeta^{-1})(1-\mu_k)^2(\eta^{\ps}_k)^2\norm{g^{\ps}_k}^2 .
}\]
By definition $\eta^{\ps}_k=\gamma_k/\max\{1,\norm{g^{\ps}_k}\}$, so
$\eta^{\ps}_k\norm{g^{\ps}_k}\le\gamma_k$ and hence
$(\eta^{\ps}_k)^2\norm{g^{\ps}_k}^2\le\gamma_k^2$; with $(1-\mu_k)^2\le 1$,
\begin{equation}\label{eq:p6}\textcolor{red}{
\norm{d_k}^2\le(1+\zeta)\rho_{\cvx}^2\mu_k^2 T_k+(1+\zeta^{-1})\gamma_k^2 .
}\end{equation}
\leavevmode\textcolor{red}{Choose $\zeta:=(2-\rho_{\cvx})/(2\rho_{\cvx})>0$. Then
$(1+\zeta)\rho_{\cvx}=(2+\rho_{\cvx})/2$, so
$(1+\zeta)\rho_{\cvx}^2=\rho_{\cvx}(2+\rho_{\cvx})/2$, and
$1+\zeta^{-1}=(2+\rho_{\cvx})/(2-\rho_{\cvx})=C_0$}. Combining
\eqref{eq:p1}, \eqref{eq:p2}, \eqref{eq:p4} and~\eqref{eq:p6},
\[\textcolor{red}{
\norm{x^{k+1}-x^\star}^2\le\norm{x^k-x^\star}^2+\bigl[-2\mu_k\rho_{\cvx}+(1+\zeta)\rho_{\cvx}^2\mu_k^2\bigr]T_k-2(1-\mu_k)\eta^{\ps}_k R^{x^\star}_k+C_0\gamma_k^2 .
}\]
The bracket equals
$\mu_k\rho_{\cvx}\bigl(-2+\tfrac{2+\rho_{\cvx}}{2}\mu_k\bigr)$. Since $\rho_{\cvx}<2$, for
all $\mu_k\in[0,1]$ the factor $-2+\tfrac{2+\rho_{\cvx}}{2}\mu_k$ is at most
$-(2-\rho_{\cvx})/2<0$ (its maximum over $[0,1]$ is at $\mu_k=1$); multiplying by
$\mu_k\rho_{\cvx}\ge0$ gives the bracket $\le -c_0\mu_k$. This yields~\eqref{eq:fejer}.
\end{proof}

\begin{lemma}\label{lem:bounded}
The sequence $\{x^k\}$ is bounded. Moreover, for each fixed $x^\star\in\Omega$:
\textup{(i)} $\{\norm{x^k-x^\star}^2\}$ converges in $[0,\infty)$;
\textup{(ii)} $\sum_{k=0}^\infty\mu_k T_k<\infty$;
\textup{(iii)} $\sum_{k=0}^\infty(1-\mu_k)\eta^{\ps}_k R^{x^\star}_k<\infty$.
\end{lemma}
\begin{proof}
Fix $x^\star\in\Omega$ and set $a_k:=\norm{x^k-x^\star}^2$. By Lemma~\ref{lem:sign-ps},
$R^{x^\star}_k\ge0$, and $\mu_k,T_k,(1-\mu_k)\eta^{\ps}_k\ge0$, so~\eqref{eq:fejer} gives
$a_{k+1}\le a_k+C_0\gamma_k^2$. Since $\sum_k C_0\gamma_k^2<\infty$ by (A3)(iii),
Lemma~\ref{lem:seq} yields (i), and boundedness of $\{x^k\}$ follows. Summing
\eqref{eq:fejer} from $0$ to $N-1$,
\[
a_N+\sum_{k=0}^{N-1}\bigl(c_0\mu_k T_k+2(1-\mu_k)\eta^{\ps}_k R^{x^\star}_k\bigr)\le a_0+C_0\sum_{k=0}^{N-1}\gamma_k^2 .
\]
The right side is bounded uniformly in $N$; letting $N\to\infty$ and using $c_0>0$ with
nonnegative terms gives (ii) and (iii).
\end{proof}

\begin{lemma}\label{lem:jcvx0}
$\displaystyle\lim_{k\to\infty}J^{\cvx}_k=0$.
\end{lemma}
\begin{proof}
If $I^{\cvx}=\emptyset$ then $J^{\cvx}_k=0$ for all $k$. Assume $I^{\cvx}\ne\emptyset$, \leavevmode\textcolor{red}{fix
$x^\star\in\Omega$}, and put $B:=\overline{\conv}\{x^k:k\ge0\}$, compact by Lemma~\ref{lem:bounded} \leavevmode\textcolor{red}{and contained
in $K\subset U$}. By (A2) and
compactness,
\[
L:=\max_{i\in I^{\cvx}}\sup_{x\in B}\norm{\nabla F_i(x)}<\infty,\qquad
V:=\max_{i\in I}\sup_{x\in B}\abs{F_i(x)-b_i}<\infty.
\]
If $V=0$ then all $v_i(x^k)=0$ and
$J^{\cvx}_k=0$; assume $V>0$. If $L=0$ then each $F_i$, $i\in I^{\cvx}$, is constant on
$B$, so $g^{\cvx}_k=0$ and, by Remark~\ref{rem:grad-cvx-sum},
$0=\inner{g^{\cvx}_k}{x^k-x^\star}\ge 2\Phi^{\cvx}_k\ge(J^{\cvx}_k)^2$, giving
$J^{\cvx}_k=0$; assume $L>0$. Set $m^{\cvx}:=\abs{I^{\cvx}}\ge1$.

Consider a step with $J^{\cvx}_k>0$ and pick $i^\star\in I^{\cvx}$ with
$v_{i^\star}(x^k)=J^{\cvx}_k$. Since $v_{i^\star}(x^k)\ge J^{\cvx}_k-\delta$, we have
$i^\star\in K^{\cvx}_k$, whence
\begin{equation}\label{eq:p8}
\Phi^{\cvx}_k\ge\tfrac12 v_{i^\star}(x^k)^2=\tfrac12(J^{\cvx}_k)^2>0 .
\end{equation}
Remark~\ref{rem:grad-cvx-sum} gives
$\inner{g^{\cvx}_k}{x^k-x^\star}\ge2\Phi^{\cvx}_k>0$, so $g^{\cvx}_k\ne0$; with
Lemma~\ref{lem:bound-cvx}, $\norm{g^{\cvx}_k}^2\le2m^{\cvx}L^2\Phi^{\cvx}_k$, and with
\eqref{eq:p8},
\begin{equation}\label{eq:p9}
T_k=\frac{(2\Phi^{\cvx}_k)^2}{\norm{g^{\cvx}_k}^2}\ge\frac{2\Phi^{\cvx}_k}{m^{\cvx}L^2}\ge\frac{(J^{\cvx}_k)^2}{m^{\cvx}L^2}.
\end{equation}
For a lower bound on $\mu_k$: for all $i\in I$ and $x^k\in B$,
$v_i(x^k)\le\abs{F_i(x^k)-b_i}\le V$, so $J^{\cvx}_k\le V$ and $J^{\ps}_k\le V$; when
$J^{\cvx}_k>0$,
\begin{equation}\label{eq:p10}
\mu_k=\frac{J^{\cvx}_k}{J^{\cvx}_k+J^{\ps}_k}\ge\frac{J^{\cvx}_k}{2V}.
\end{equation}
Combining~\eqref{eq:p9} and~\eqref{eq:p10}, for $J^{\cvx}_k>0$,
$\mu_k T_k\ge(J^{\cvx}_k)^3/(2Vm^{\cvx}L^2)$; this also holds (as $0\ge0$) when
$J^{\cvx}_k=0$. By Lemma~\ref{lem:bounded}(ii),
\begin{equation}\label{eq:p12}
\sum_{k=0}^\infty(J^{\cvx}_k)^3\le 2Vm^{\cvx}L^2\sum_{k=0}^\infty\mu_k T_k<\infty,
\end{equation}
so $(J^{\cvx}_k)^3\to0$ and thus $J^{\cvx}_k\to0$.
\end{proof}

\begin{proposition}\label{prop:regularity}
\[\textcolor{red}{
\sum_{k=0}^{\infty}\norm{x^{k+1}-x^k}^{2}<\infty;
}\]
in particular,
$\norm{x^{k+1}-x^k}\to0$.
\end{proposition}
\begin{proof}
Since $x^k\in K$ and $x^{k+1}=P_K(y^k)$, Lemma~\ref{lem:proj-nonexp} with $y=y^k$, $z=x^k$
gives $\norm{x^{k+1}-x^k}\le\norm{y^k-x^k}=\norm{d_k}$. By~\eqref{eq:p6} and
$\mu_k^2\le\mu_k$, \leavevmode\textcolor{red}{$\norm{d_k}^2\le(1+\zeta)\rho_{\cvx}^2\mu_k T_k+C_0\gamma_k^2$}.
Summing and using Lemma~\ref{lem:bounded}(ii) and (A3)(iii), both series are finite.
\end{proof}

\subsection{\textcolor{red}{Convergence to a feasible point}}\label{ssec:strong}
\leavevmode\textcolor{red}{We now prove that the whole sequence converges to a feasible point. The proof has three
stages. Stage~1 uses Lemma~\ref{lem:jcvx0} to show that every cluster point satisfies the
convex constraints. Stage~2 shows, by contradiction and pseudoconvexity, that at least one
cluster point is feasible. Stage~3 uses the convergence of $\norm{x^k-\bar x}$ to show that
the whole sequence converges to that point.}

\begin{theorem}\label{thm:strong}
Under Assumption~\ref{ass:standing}, the sequence $\{x^k\}$ generated by
Algorithm~\ref{alg:dtap} converges to a point $x^\infty\in\Omega$.
\end{theorem}
\begin{proof}
By Lemma~\ref{lem:bounded}, $\{x^k\}$ is bounded; let
$C:=\{\bar x\in H:\exists\,k_l\to\infty,\ x^{k_l}\to\bar x\}$ be its (nonempty, compact)
cluster set. Since $K$ is closed and $x^k\in K$, $C\subseteq K$.

\emph{Stage 1 (convex \leavevmode\textcolor{red}{constraints).}} Fix $\bar x\in C$ and $x^{k_l}\to\bar x$. By (A2) each
$F_i$ is continuous, so $v_i(x^{k_l})\to v_i(\bar x)$; for $i\in I^{\cvx}$,
$v_i(x^{k_l})\le J^{\cvx}_{k_l}\to0$ by Lemma~\ref{lem:jcvx0}, so by squeezing
$v_i(\bar x)=0$, i.e.\ $F_i(\bar x)\le b_i$. Hence
\begin{equation}\label{eq:p13}\textcolor{red}{
C\subseteq\bigcap_{i\in I^{\cvx}}\{x\in K:F_i(x)\le b_i\}.
}\end{equation}
If $I^{\ps}=\emptyset$, then~\eqref{eq:p13} gives \leavevmode\textcolor{red}{$C\subseteq\Omega$},
so $C\cap\Omega=C\ne\emptyset$.
Assume henceforth $I^{\ps}\ne\emptyset$.

\emph{\leavevmode\textcolor{red}{Stage 2 (a feasible cluster point).}} Suppose, for contradiction, that
$C\cap\Omega=\emptyset$. By~\eqref{eq:p13} and the partition $I=I^{\cvx}\cup I^{\ps}$, for
each $\bar x\in C$ there is $j\in I^{\ps}$ with $F_j(\bar x)>b_j$, i.e.\
$J^{\ps}(\bar x)>0$. As a maximum of finitely many continuous functions, $J^{\ps}$ is
continuous and positive on the compact set $C$, so $\alpha:=\min_{x\in C}J^{\ps}(x)>0$.

\smallskip
\emph{Claim 1 (uniform lower bound for $J^{\ps}_k$).} There exist $\tau_0>0$ and $k_0$
with $J^{\ps}_k\ge\tau_0$ for all $k\ge k_0$. Indeed, put $\tau_0:=\alpha/2$. If the claim
failed there would be $k_l\to\infty$ with $J^{\ps}_{k_l}<\tau_0$; extracting (same
notation) $x^{k_l}\to\tilde x\in C$, continuity gives $J^{\ps}(\tilde x)\le\tau_0<\alpha$,
contradicting the definition of $\alpha$.

\smallskip
\emph{Claim 2 ($\mu_k\to0$).} For $k\ge k_0$,
$\mu_k=J^{\cvx}_k/(J^{\cvx}_k+J^{\ps}_k)\le J^{\cvx}_k/\tau_0\to0$ by
Lemma~\ref{lem:jcvx0}. Hence there is $k_1\ge k_0$ with $1-\mu_k\ge1/2$ for $k\ge k_1$.

\smallskip
\emph{Claim 3 (uniform lower bound for $R^{x^\star}_k$).} Fix $x^\star\in\Omega$. There
exist $\xi_0>0$ and $k_2\ge k_1$ with $R^{x^\star}_k\ge\xi_0$ for all $k\ge k_2$. Suppose
not: there is $k_l\to\infty$ with $R^{x^\star}_{k_l}\to0$. As $\{x^{k_l}\}$ is bounded,
extract $x^{k_l}\to\bar x\in C$; since $I^{\ps}$ has finitely many subsets, extract
further so that $K^{\ps}_{k_l}\equiv K^\star$ is constant. By assumption
$J^{\ps}(\bar x)>0$; choose $j^\star\in\argmax_{j\in I^{\ps}}v_j(\bar x)$, so
$v_{j^\star}(\bar x)=J^{\ps}(\bar x)>0$. By continuity $v_{j^\star}(x^{k_l})\to
v_{j^\star}(\bar x)$ and $J^{\ps}_{k_l}\to J^{\ps}(\bar x)$, so
$J^{\ps}_{k_l}-v_{j^\star}(x^{k_l})\to0$; with $\delta>0$ from (A3)(i), for large $l$ we
get $v_{j^\star}(x^{k_l})\ge J^{\ps}_{k_l}-\delta$, hence $j^\star\in K^{\ps}_{k_l}=K^\star$.
Each summand of
$R^{x^\star}_{k_l}=\sum_{j\in K^\star}v_j(x^{k_l})\inner{\nabla F_j(x^{k_l})}{x^{k_l}-x^\star}$
is nonnegative by Lemma~\ref{lem:sign-ps}, so keeping only $j^\star$,
\[
R^{x^\star}_{k_l}\ge v_{j^\star}(x^{k_l})\inner{\nabla F_{j^\star}(x^{k_l})}{x^{k_l}-x^\star}.
\]
Letting $l\to\infty$ and using continuity of $v_{j^\star}$ and $\nabla F_{j^\star}$,
\[
\liminf_{l\to\infty}R^{x^\star}_{k_l}\ge v_{j^\star}(\bar x)\inner{\nabla F_{j^\star}(\bar x)}{\bar x-x^\star}=:\Lambda .
\]
Since $j^\star\in I^{\ps}$ and $x^\star\in\Omega$, $F_{j^\star}(x^\star)\le b_{j^\star}<F_{j^\star}(\bar x)$;
definition~\eqref{eq:contrapositive} with $x=\bar x$, $y=x^\star$ gives
$\inner{\nabla F_{j^\star}(\bar x)}{\bar x-x^\star}>0$, and with $v_{j^\star}(\bar x)>0$ we
get $\Lambda>0$, contradicting $R^{x^\star}_{k_l}\to0$. This proves Claim~3.

\smallskip
\emph{Reaching the contradiction.} By Lemma~\ref{lem:bound-ps} on
$B:=\overline{\conv}\{x^k\}$, \leavevmode\textcolor{red}{$\norm{g^{\ps}_k}\le\abs{I^{\ps}}V_B L'_B$ with
$L'_B:=\max_{j\in I^{\ps}}\sup_B\norm{\nabla F_j}<\infty$}. Put
\leavevmode\textcolor{red}{$G':=\max\{1,\abs{I^{\ps}}V_B L'_B\}$}; then $\max\{1,\norm{g^{\ps}_k}\}\le G'$ for all $k$,
so $\eta^{\ps}_k=\gamma_k/\max\{1,\norm{g^{\ps}_k}\}\ge\gamma_k/G'$. For $k\ge k_2$,
combining with $1-\mu_k\ge1/2$ and $R^{x^\star}_k\ge\xi_0$,
\[
\sum_{k=0}^\infty(1-\mu_k)\eta^{\ps}_k R^{x^\star}_k\ge\sum_{k\ge k_2}\frac12\cdot\frac{\gamma_k}{G'}\cdot\xi_0=\frac{\xi_0}{2G'}\sum_{k\ge k_2}\gamma_k=+\infty
\]
by (A3)(iii), contradicting Lemma~\ref{lem:bounded}(iii). Hence $C\cap\Omega\ne\emptyset$.

\emph{Stage 3 (whole-sequence convergence).} Fix $\bar x\in C\cap\Omega$.
Lemma~\leavevmode\textcolor{red}{\ref{lem:bounded}(i), applied to $x^\star=\bar x$, ensures that
$\lim_{k\to\infty}\norm{x^k-\bar x}^2$ exists. Since $\bar x\in C$, a subsequence
$x^{k_l}$ converges to $\bar x$. Therefore}
\[\textcolor{red}{
\lim_{k\to\infty}\norm{x^k-\bar x}^2
=\lim_{l\to\infty}\norm{x^{k_l}-\bar x}^2=0.
}\]
\leavevmode\textcolor{red}{Hence the whole sequence converges to $\bar x$, and we set $x^\infty:=\bar x\in\Omega$.
The argument does not require $\Omega$ to be a singleton: one feasible cluster point and
the convergence of its distance suffice.}
\end{proof}

\begin{remark}\label{rem:delta-role}
The parameter $\delta>0$ is not necessary for convergence. When $\delta=0$ one argues by
subsequence: along $x^{k_l}\to\bar x\in C$, choose
$j_l\in\argmax_{j\in I^{\ps}}v_j(x^{k_l})$; since $I^{\ps}$ is finite, extract
$j_l\equiv j^\star$, and $v_{j^\star}(x^{k_l})=J^{\ps}(x^{k_l})\to J^{\ps}(\bar x)>0$, so
$j^\star$ attains the maximal violation at $\bar x$ and the pseudoconvexity argument
goes through. The role of $\delta>0$ is to simplify: a maximizer at $\bar x$ automatically
lies in $K^{\ps}_{k_l}$ for large $l$, avoiding a `finitely many indices' extraction.
Larger $\delta$ enlarges \leavevmode\textcolor{red}{$K^{\cvx}_k$ and $K^{\ps}_k$, which may improve} numerical
stability when several $v_i(x^k)$ are nearly equal.
\end{remark}

\begin{remark}\label{rem:rate}
From~\eqref{eq:p12}, $\sum_k(J^{\cvx}_k)^3<\infty$, whence
\leavevmode\textcolor{red}{$\min_{0\le k<N}J^{\cvx}_k=O(N^{-1/3})$. This bounds the \emph{best} convex-track violation
over the first $N$ iterates. It is not a last-iterate bound, a rate for the full residual
$J(x^k)$, or a rate for the distance to a feasible point. For these latter quantities,
under (A1)--(A3) alone we assert only qualitative convergence.}
\end{remark}

\subsection{Degenerate constraint structures}\label{ssec:degenerate}
When the constraint structure degenerates (only a convex track or only a pseudoconvex
track), \dtap{} and its analysis simplify, yielding two corollaries of
Theorem~\ref{thm:strong}.

\begin{corollary}\label{cor:cvx}
If $I^{\ps}=\emptyset$, then $J^{\ps}_k\equiv0$, $g^{\ps}_k\equiv0$, and $\mu_k\equiv1$
whenever $J^{\cvx}_k>0$, so Algorithm~\ref{alg:dtap} reduces to
$x^{k+1}=P_K(x^k-\eta^{\cvx}_k g^{\cvx}_k)$. The sequence converges to a point
\leavevmode\textcolor{red}{$x^\infty\in\Omega=\{x\in K:F_i(x)\le b_i\text{ for all }i\in I^{\cvx}\}$}.
\end{corollary}
\begin{proof}
\leavevmode\textcolor{red}{Here $d_k=\eta^{\cvx}_k g^{\cvx}_k$, and~\eqref{eq:p1} directly gives
$\norm{x^{k+1}-x^\star}^2\le\norm{x^k-x^\star}^2-2\eta^{\cvx}_k\inner{g^{\cvx}_k}{x^k-x^\star}+(\eta^{\cvx}_k)^2\norm{g^{\cvx}_k}^2$}.
If $\norm{g^{\cvx}_k}=0$ then $\eta^{\cvx}_k=0$, $T_k=0$, and the bound is
$\norm{x^{k+1}-x^\star}^2\le\norm{x^k-x^\star}^2$. Otherwise Lemma~\ref{lem:grad-cvx} gives
$\inner{g^{\cvx}_k}{x^k-x^\star}\ge2\Phi^{\cvx}_k$, and substituting
$\eta^{\cvx}_k=2\rho_{\cvx}\Phi^{\cvx}_k/\norm{g^{\cvx}_k}^2$ together with
$(\eta^{\cvx}_k)^2\norm{g^{\cvx}_k}^2=\eta^{\cvx}_k\cdot2\rho_{\cvx}\Phi^{\cvx}_k$ yields, after
simplification,
$2\eta^{\cvx}_k\inner{g^{\cvx}_k}{x^k-x^\star}-(\eta^{\cvx}_k)^2\norm{g^{\cvx}_k}^2\ge\rho_{\cvx}(2-\rho_{\cvx})(2\Phi^{\cvx}_k)^2/\norm{g^{\cvx}_k}^2=2c_0T_k$.
Hence
\begin{equation}\label{eq:p15}
\norm{x^{k+1}-x^\star}^2\le\norm{x^k-x^\star}^2-2c_0T_k ,
\end{equation}
valid for all $k$. Thus $\{\norm{x^k-x^\star}^2\}$ is nonincreasing and convergent, and
summing gives $\sum_k T_k<\infty$. As in Lemma~\ref{lem:jcvx0}, when $J^{\cvx}_k>0$ one has
$T_k\ge(J^{\cvx}_k)^2/(m^{\cvx}L^2)$ (note the exponent $2$ here, versus $3$
in~\eqref{eq:p12}, since no $\mu_k$ factor is present), so $\sum_k(J^{\cvx}_k)^2<\infty$ and
$J^{\cvx}_k\to0$. Finally, $I^{\ps}=\emptyset$ makes Stage~2 of
Theorem~\ref{thm:strong} vacuous (indeed $C\subseteq\Omega$ already
from~\eqref{eq:p13}), and Stage~3 gives $x^k\to x^\infty\in\Omega$.
\end{proof}

\begin{corollary}\label{cor:ps}
If $I^{\cvx}=\emptyset$, then $J^{\cvx}_k\equiv0$, $g^{\cvx}_k\equiv0$, and $\mu_k\equiv0$
whenever $J^{\ps}_k>0$, so Algorithm~\ref{alg:dtap} reduces to
$x^{k+1}=P_K(x^k-\eta^{\ps}_k g^{\ps}_k)$ with $\eta^{\ps}_k=\gamma_k/\max\{1,\norm{g^{\ps}_k}\}$.
The sequence converges to a point $x^\infty\in\Omega$.
\end{corollary}
\begin{proof}
Now $d_k=\eta^{\ps}_k g^{\ps}_k$ with $\norm{d_k}^2\le\gamma_k^2$, so~\eqref{eq:p1} gives
$\norm{x^{k+1}-x^\star}^2\le\norm{x^k-x^\star}^2-2\eta^{\ps}_k R^{x^\star}_k+\gamma_k^2$,
which is~\eqref{eq:fejer} in this degenerate case ($T_k=0$, $1-\mu_k=1$). Since
$R^{x^\star}_k\ge0$ (Lemma~\ref{lem:sign-ps}), $a_{k+1}\le a_k+\gamma_k^2$, and
Lemma~\ref{lem:seq} gives convergence of $\{\norm{x^k-x^\star}^2\}$ and boundedness;
summation gives $\sum_k\eta^{\ps}_k R^{x^\star}_k<\infty$. Stage~1 of
Theorem~\ref{thm:strong} gives $C\subseteq K$ (the convex part being vacuous); Stage~2 with
$\mu_k=0$ (so $1-\mu_k=1$) produces the same contradiction $\sum\eta^{\ps}_k R^{x^\star}_k=+\infty$,
so $C\cap\Omega\ne\emptyset$; Stage~3 gives $x^k\to x^\infty\in\Omega$.
\end{proof}


\section{Numerical experiments}\label{sec:experiments}
\leavevmode\textcolor{red}{We separate predictive performance from the numerical solution of the feasibility system. The classification comparison uses eight datasets. The solver diagnostic additionally includes Movement Libras and controlled nonconvex pseudoconvex instances.}

\subsection{\textcolor{red}{Feature-selection model}}\label{sec:fs}
\leavevmode\textcolor{red}{We apply the normal-box NSFP to feature selection by combining criteria from the fractional-programming model of Wang et al.~\citep{wang2021fractional} and the fuzzy-rough model AFRSFCV~\citep{wang2026afrsfcv}. Let $\mathcal{X}=\{x_1,\ldots,x_n\}$ be the training samples, $A=\{a_1,\ldots,a_p\}$ the features, $X\in\R^{n\times p}$ the data matrix, and $y\in\{1,\ldots,r\}^n$ the class labels. Here $n$ is the number of samples; subscripts index features and superscripts index iterations. The feature weights belong to the compact convex simplex $K=\Delta_{p-1}=\{w\in\R^p:w\ge0,\ \mathbf1_p^{\top}w=1\}$. All supervised quantities below are estimated from the training data.}

\subsubsection{\textcolor{red}{Fractional criterion and N-Fscore}}\label{ssec:wang}\label{sssec:wang-inputs}\label{sssec:nfscore}
The supervised similarity matrix $S\in\R^{p\times p}$ has entries
\begin{equation}\label{eq:simMI}\textcolor{red}{
s_{ij}=\max\left\{0,\frac{I(Z_i;Z_j;C)}{H(Z_i)+H(Z_j)}\right\},\qquad i,j=1,\ldots,p,
}\end{equation}
\leavevmode\textcolor{red}{where $Z_i$ is the discretized feature $a_i$, $C$ is the class variable, and $I(Z_i;Z_j;C)=I(Z_i;C)+I(Z_j;C)-I((Z_i,Z_j);C)$ is their interaction information.} Entropy and \leavevmode\textcolor{red}{mutual information} are estimated from equal-width histograms with \leavevmode\textcolor{red}{ten} bins per feature~\citep{brown2012}. \leavevmode\textcolor{red}{We set $s_{ij}=0$ if the denominator is zero. The redundancy matrix is}
\begin{equation}\label{eq:matrixQ}\textcolor{red}{
Q_R=S+\delta_QI_p,\qquad \delta_Q\ge-\min\{0,\lambda_{\min}(S)\},
}\end{equation}
\leavevmode\textcolor{red}{so $Q_R\succeq0$. The implementation uses $\delta_Q=\max\{0,-\lambda_{\min}(S)\}+10^{-10}$. Let $n_t$, $\bar a_{kt}$, and $\sigma_{kt}^2$ be the class size, mean, and population variance of feature $a_k$ in class $t$, and let $\bar a_k$ be its overall mean. The Fisher scores are}
\begin{equation}\label{eq:rhoFisher}\textcolor{red}{
\rho_k=\frac{\sum_{t=1}^r n_t(\bar a_{kt}-\bar a_k)^2}{\sum_{t=1}^r n_t\sigma_{kt}^2+\varepsilon_\rho},\qquad \varepsilon_\rho=10^{-12}.
}\end{equation}
\leavevmode\textcolor{red}{These scores are nonnegative. The denominator regularization prevents division by zero but does not ensure $\rho_k>0$, since the numerator can vanish.}

\leavevmode\textcolor{red}{The fractional program underlying N-Fscore is}
\begin{equation}\label{eq:wangOrig}\textcolor{red}{
\min_{w\in K}\ \frac{w^{\top}Q_Rw}{\rho^{\top}w},
}\end{equation}
\leavevmode\textcolor{red}{assuming the denominator is positive on $K$. The numerator measures redundancy and the denominator relevance. Our N-Fscore baseline uses the projection-neural-network stage of~\citep{wang2021fractional}, with the Fisher scores in~\eqref{eq:rhoFisher}. It then applies the common selection rule~\eqref{eq:thetaRule}, instead of the source method's $k$-winners-take-all stage. Its dynamics and explicit-Euler implementation are given in Section~\ref{sssec:baselines}}.

\subsubsection{\textcolor{red}{Fuzzy-rough criterion and AFRSFCV}}\label{ssec:afrsfcv}\label{sssec:afrsfcv-inputs}\label{sssec:afrsfcv-model}
\leavevmode\textcolor{red}{For each class, write $D_t=\{x_i:y_i=t\}$ and $D_t(x_j)=\mathbf1_{\{y_j=t\}}$. AFRSFCV incorporates class variance into a fuzzy relation. With $a_k(x_i)=X_{ik}$, the implemented relation is}
\begin{equation}\label{eq:simRel}\textcolor{red}{
R_k(x_i,x_j)=\exp\!\left(-\frac{[a_k(x_i)-a_k(x_j)]^2}{2\widetilde\sigma_{k,y_i}^2}\right),
\qquad \widetilde\sigma_{kt}^2=\max\{\sigma_{kt}^2,10^{-10}\}.
}\end{equation}
\leavevmode\textcolor{red}{The variance floor makes the formula defined when a feature is constant within a class. The fuzzy lower approximation is}
\begin{equation}\label{eq:lower}\textcolor{red}{
\underline{R_k}(D_t)(x_i)=\min_{x_j\in\mathcal{X}}\max\{1-R_k(x_i,x_j),D_t(x_j)\}.
}\end{equation}
\leavevmode\textcolor{red}{Its values form the matrix $\underline{R_k}(D)\in[0,1]^{n\times r}$. Let $Y_{it}=\mathbf1_{\{y_i=t\}}$ be the label matrix and define the decision matrix by}
\begin{equation}\label{eq:Pmat}\textcolor{red}{
P_{ik}=\mathbf1_{\{\argmax_{1\le t\le r}\underline{R_k}(D_t)(x_i)=y_i\}}.
}\end{equation}
\leavevmode\textcolor{red}{Ties are resolved by the smallest class index. Thus $P_{ik}=1$ precisely when feature $a_k$ predicts sample $x_i$ correctly under this rule.}

\leavevmode\textcolor{red}{The AFRSFCV baseline~\citep{wang2026afrsfcv} minimizes}
\begin{equation}\label{eq:afrsfcvLoss}
\begin{split}
\min_{\tau\in\R^p}\ L(\tau)
&=\frac12\left\|Y-\sum_{k=1}^{p}w_k(\tau)\underline{R_k}(D)\right\|_F^2\\
&\quad+\frac{\lambda_1}{2}\|\mathbf1_n-Pw(\tau)\|^2+\frac{\lambda_2}{2}\|w(\tau)\|^2,
\end{split}
\end{equation}
\leavevmode\textcolor{red}{where $w_k(\tau)=e^{\tau_k}/\sum_{l=1}^{p}e^{\tau_l}$ ensures $w(\tau)\in K$. The three terms are the label-approximation error, the weight of misclassifying features per sample, and a Tikhonov term. The baseline uses gradient descent for a fixed iteration budget, followed by~\eqref{eq:thetaRule}}. The gradient and update are given in Section~\ref{sssec:baselines}.

\subsubsection{\textcolor{red}{Feasibility and feature selection}}\label{ssec:reduction}\label{ssec:selection}
\leavevmode\textcolor{red}{We impose separate bounds on three criteria,}
\begin{align}\color{red}
F_1(w)&\color{red}=\left\|Y-\sum_{k=1}^{p}w_k\underline{R_k}(D)\right\|_F^2,\label{eq:fsF1}\\\color{red}
F_2(w)&\color{red}=\|\mathbf1_n-Pw\|^2,\label{eq:fsF2}\\\color{red}
F_3(w)&\color{red}=\frac{w^{\top}Q_Rw}{\rho^{\top}w},\qquad U_\rho=\{w:\rho^{\top}w>0\}.\label{eq:fsF3}
\end{align}
Here $F_1,F_2$ are the first two terms of~\eqref{eq:afrsfcvLoss}, \leavevmode\textcolor{red}{without their coefficients,} and $F_3$ is the objective \leavevmode\textcolor{red}{in}~\eqref{eq:wangOrig}. \leavevmode\textcolor{red}{For thresholds $b_i>0$, the resulting problem is}
\begin{equation}\label{eq:nsfp-fs}
\text{find }w^\star\in\Omega,\qquad \Omega=\{w\in K:F_i(w)\le b_i,\ i=1,2,3\}.
\end{equation}
\leavevmode\textcolor{red}{This formulation asks for simultaneous acceptability of the three criteria. It omits the Tikhonov term and does not preserve either source optimization problem or its minimizers.}

\begin{proposition}\label{prop:fs-reg}
\leavevmode\textcolor{red}{If $Q_R\succeq0$ and $\min_k\rho_k>0$, all three functions in~\eqref{eq:fsF1}--\eqref{eq:fsF3} are smooth and convex on an open convex neighbourhood of $K$. Consequently,~\eqref{eq:nsfp-fs} is an all-convex normal-box NSFP. Hence~\eqref{eq:nsfp-fs} satisfies (A1) and (A2) of Assumption~\ref{ass:standing}, with $U=U_\rho$, whenever $\Omega\ne\varnothing$}.
\end{proposition}
\begin{proof}
\leavevmode\textcolor{red}{The first two functions are squared norms of affine maps. For $w\in U_\rho$ and any direction $h$, put $t=\rho^{\top}w$ and $a=\rho^{\top}h$. Direct differentiation gives}
\begin{equation}\label{eq:fsF3-curvature}\textcolor{red}{
D^2F_3(w)[h,h]=\frac2t\left(h-\frac atw\right)^{\!\top}Q_R\left(h-\frac atw\right)\ge0.
}\end{equation}
\leavevmode\textcolor{red}{Thus $F_3$ is convex and smooth on the open half-space $U_\rho$. Since $\rho^{\top}w\ge\min_k\rho_k>0$ for $w\in K$, this half-space contains $K$}.
\end{proof}

\leavevmode\textcolor{red}{The experiments use $I^{\cvx}=\{1,2\}$, $I^{\ps}=\{3\}$, which is admissible since convex functions are pseudoconvex; DTAP then uses only sign information for $F_3$, so no genuinely nonconvex constraint is tested. The all-convex assignment $I^{\cvx}=\{1,2,3\}$ is also admissible, and Theorem~\ref{thm:strong} applies to both. Proposition~\ref{prop:fs-reg} needs $\min_k\rho_k>0$; on all instances used here the smallest score is $4.77\times10^{-5}$. Hence the safeguard $\max\{\rho^\top w,10^{-15}\}$ in the implementation is never active, and the implemented $F_3$ is~\eqref{eq:fsF3}}.

\leavevmode\textcolor{red}{The thresholds are calibrated on the training data: $b_i$ is the $q_i$-th percentile of $F_i$ over $3\,000$ Dirichlet points (seed $42$) and the $p$ vertices of $K$, where a fixed rule chooses $q_i$ from the sample size, the number of features and classes, and the class imbalance (supplementary metadata). All solvers share these thresholds. Calibrating each $b_i$ separately does not guarantee $\Omega\ne\varnothing$; Section~\ref{ssec:feasibility-audit} checks feasibility numerically. Prescribed thresholds would be used as given.}

\leavevmode\textcolor{red}{For every method, sort the computed weights as $w^\star_{(1)}\ge\cdots\ge w^\star_{(p)}$ and retain the first $s$ features, where}
\begin{equation}\label{eq:thetaRule}\textcolor{red}{
s=\min\left\{s'\in\{1,\ldots,p\}:\sum_{j=1}^{s'}w^\star_{(j)}\ge\theta\right\},\qquad 0<\theta<1.
}\end{equation}
\leavevmode\textcolor{red}{This cumulative-weight rule does not enforce a prescribed feature budget. Truncating the weights does not preserve the feasibility guarantee for the full vector.}

\begin{remark}\label{rem:fs-cost}
\leavevmode\textcolor{red}{With the precomputed matrices, evaluating $F_1,\nabla F_1$ costs $O(nrp)$, evaluating $F_2,\nabla F_2$ costs $O(np)$, and evaluating $F_3,\nabla F_3$ costs $O(p^2)$. Projection onto $K=\Delta_{p-1}$ by sorting costs $O(p\log p)$ operations~\citep{condat2016simplex}. A dense iteration therefore costs $O(nrp+p^2)$ operations and requires $O(nrp+p^2)$ storage, including $Q_R$. Preprocessing and classifier training are separate costs.}
\end{remark}

\subsection{Experimental setup}\label{ssec:setup}
\leavevmode\textcolor{red}{The main classification study uses the eight UCI datasets~\citep{kelly2023uci} in Table~\ref{tab:datasets}. The DTAP runs reported here start at the uniform weight vector and use $\rho_{\cvx}=1.98$, $\delta=10^{-3}$, $\gamma_k=(k+0.5)^{-0.55}$, tolerance $10^{-6}$, and a budget of $1000$ iterations. SG-const uses the step $\nu_k=10^{-3}$ and SG-dyn the sampled modulus $\widehat L$ (Section~\ref{sssec:baselines}); tuned moduli are considered only in Section~\ref{ssec:feasibility-audit}. N-Fscore uses $\varepsilon_{\mathrm{PNN}}=10^{-3}$ and $1000$ Euler grid points on $[0,0.2]$. AFRSFCV uses learning rate $10^{-3}$, $\lambda_1=1$, $\lambda_2=0.5$, and $1000$ steps.} The shared cumulative threshold is $\theta=0.9$.

\begin{table}[htbp]
\centering
\caption{\textcolor{red}{Datasets used for the main classification comparison. The additional Movement Libras instance is used only in the solver diagnostic.}}\label{tab:datasets}
\begin{tabular}{lrrrl}
\toprule
Dataset      & Samples $n$ & Features $p$ & Classes $r$ & Balance    \\
\midrule
Glass        & $214$       & $9$          & $6$         & Imbalanced \\
Diabetes     & $768$       & $8$          & $2$         & Imbalanced \\
Segmentation & $2\,310$    & $18$         & $7$         & Balanced   \\
Abalone      & $4\,177$    & $8$          & $3$         & Balanced   \\
Satimage     & $6\,435$    & $36$         & $6$         & Imbalanced \\
Yeast        & $1\,484$    & $7$          & $10$        & Imbalanced \\
Spambase     & $4\,601$    & $57$         & $2$         & Imbalanced \\
Pendigits    & $10\,992$   & $16$         & $10$        & Balanced   \\
\bottomrule
\end{tabular}
\end{table}
\subsubsection{\textcolor{red}{Baseline solvers}}\label{sssec:baselines}
\leavevmode\textcolor{red}{The following paragraphs specify the baseline solvers: N-Fscore and AFRSFCV for the models of Sections~\ref{sssec:nfscore} and~\ref{sssec:afrsfcv-model}, and the subgradient methods SG-const and SG-dyn.}

\paragraph*{N-Fscore: the projection neural network}
To solve the fractional program~\eqref{eq:wangOrig}, Wang et al.~\citep{wang2021fractional} proposed a neurodynamic method coupling a projection neural network (PNN) with a $k$-winners-take-all network. The PNN is the differential system
\begin{equation}\label{eq:pnn}\textcolor{red}{
\begin{split}\varepsilon_{\mathrm{PNN}}\frac{du}{dt}={}&-\Pi g(u)+\tfrac1p\mathbf{1}_p\\&-(I-\Pi)\bigl(u-g(u)+\nabla f((I-\Pi)g(u)+\tfrac1p\mathbf{1}_p)\bigr).\end{split}
}\end{equation}
with output $w=g(u)$, where $g(u)=\max\{0,u\}$ is the componentwise ReLU, \leavevmode\textcolor{red}{$\Pi=\tfrac1p\mathbf{1}_p\mathbf{1}_p^{\top}$}, $\nabla f$ is the gradient of the objective in~\eqref{eq:wangOrig}, and $\varepsilon_{\mathrm{PNN}}>0$ is a time constant. \leavevmode\textcolor{red}{We integrate~\eqref{eq:pnn} by the explicit Euler method on a uniform grid of $N_{\mathrm{grid}}$ points in $[0,T]$, without a stopping rule, using the parameters of Section~\ref{ssec:setup}. The weight vector is the projection $P_K\bigl(g(u(T))\bigr)$ of the final network output onto $K$}. For a uniform comparison we use only the PNN part to obtain $w^\star$; the $k$-WTA stage is replaced by the shared cumulative-threshold rule of Section~\ref{ssec:selection}.

\paragraph*{AFRSFCV: gradient and update}
The softmax loss~\eqref{eq:afrsfcvLoss} is minimized over $\tau\in\R^p$ by gradient descent. Its gradient is
\begin{equation}\label{eq:gradTau}
\nabla_\tau L(\tau)=\bigl(\diag(w(\tau))-w(\tau)w(\tau)^{\top}\bigr)\nabla_w L(w(\tau)),
\end{equation}
with $\nabla_w L=-\mathbf{M}^{\top}(\vecop(Y)-\mathbf{M}w)+\lambda_1(P^{\top}Pw-P^{\top}\mathbf{1}_n)+\lambda_2 w$, where $\mathbf{M}\in\R^{nr\times p}$ has $k$-th column $\vecop(\underline{R_k}(D))$. The update $\tau^{t+1}=\tau^t-\eta\,\nabla_\tau L(\tau^t)$ runs for a fixed budget of $T$ steps; the final weight $w^\star=w(\tau^T)\in K$ feeds the shared selection rule of Section~\ref{ssec:selection}.

\paragraph*{\textcolor{red}{SG-const and SG-dyn}}
\leavevmode\textcolor{red}{These subgradient methods for quasiconvex feasibility~\citep{hu2022quasi} follow its
most-violated-constraint strategy: with $\mathcal I_k=\argmax_i v_i(w^k)$},
\[\textcolor{red}{
w^{k+1}=P_K(w^k-\nu_k\bar h_k),\qquad
\bar h_k=\frac1{|\mathcal I_k|}\sum_{i\in\mathcal I_k}\frac{\nabla F_i(w^k)}{\|\nabla F_i(w^k)\|}.
}\]
\leavevmode\textcolor{red}{SG-const uses $\nu_k=10^{-3}$; SG-dyn uses the dynamic step $\nu_k=J(w^k)/L$ (H\"older
exponent $1$), with $L=\widehat L$, the largest gradient norm at $500$ points sampled in $K$}.



\subsubsection{Evaluation protocol}\label{sssec:protocol}
\leavevmode\textcolor{red}{Each dataset uses a single stratified $80/20$ training-pool/test split
(seed $42$). Min--max scaling is fitted once on the full training pool, not
separately within its selection folds, and is then applied to the test data.
Stratified folds (seed $42$) within the pool use $N_f=10$, reduced to $7$ on Glass
and $4$ on Yeast. The fold masks use $\theta=0.9$ in~\eqref{eq:thetaRule};
the fixed strict-majority rule retains features selected in at least
$\lfloor N_f/2\rfloor+1$ folds. We denote the voted feature set by $A^\star$. The final classifiers use the full pool with these
voted features: RBF-SVM ($C=1$, \texttt{gamma=scale}) and $5$-NN, with the full
feature set as a reference. Accuracy and macro one-vs-rest AUC are reported.}

\leavevmode\textcolor{red}{The voted masks are computed once and then reused: the evaluation stage trains
only the downstream classifiers on these fixed masks, with SVM probability-calibration
seed $42$; identical masks share the same predictions. The test data are not used to select features or to choose any parameter or vote threshold.
The training/test split is fixed before the selections are computed, and the rule that chooses
the threshold levels uses only the sample size, the number of features, the number of classes
and the class sizes of the whole dataset, not any feature values. Because a single split is
used, this is an exploratory evaluation, not an end-to-end or nested-validation study.
For the fold-selection matrix $Z\in\{0,1\}^{N_f\times p}$, the Nogueira index~\citep{nogueira2018}
is computed as}
\[\textcolor{red}{
\Psi=1-\frac{p^{-1}\sum_{j=1}^{p}s_j^2}{\bar q(1-\bar q)},\quad
s_j^2=\frac{1}{N_f-1}\sum_{b=1}^{N_f}(Z_{bj}-\bar Z_j)^2,\quad
\bar q=\frac{1}{N_f p}\sum_{b,j}Z_{bj}.
}\]
\leavevmode\textcolor{red}{Here $\bar Z_j=N_f^{-1}\sum_b Z_{bj}$, and $s_j^2$ is the unbiased variance over the $N_f$
folds.
Since the training folds overlap, stability is reported descriptively.
For each classifier metric, we use a two-sided signed-rank test across the
eight datasets, omitting zero differences, assigning midranks to ties, and
enumerating all signs of the nonzero differences; all-zero pairs receive the $p$-value $1$.
Holm adjustment covers the four DTAP--baseline comparisons separately for
each metric. The Friedman test compares the five selection methods.}

\subsection{Classification results}\label{ssec:results}
\leavevmode\textcolor{red}{Tables~\ref{tab:nfeat}--\ref{tab:perf-summary} report selected-set size,
stability, and predictive performance. DTAP retains $84.2\%$ of the features
on average, including $33/36$ on Satimage and $49/57$ on Spambase.
Its mean stability is $0.786$, compared with N-Fscore's $0.822$.
Thus DTAP reduces the number of features moderately; it does not select sparser or more
stable sets than the base models. AFRSFCV selects much smaller sets on several datasets,
with lower mean accuracy under the fixed configuration.}

\leavevmode\textcolor{red}{DTAP's mean accuracy is within $0.44$ percentage points of the full-feature
reference for both classifiers. It exceeds N-Fscore and AFRSFCV by $3.12$
and $4.83$ points on SVM accuracy, while retaining more features.
This is a size--prediction tradeoff, not a matched-budget comparison.
Both SG variants have exactly the same voted masks as DTAP on seven datasets;
only Spambase differs ($48$ versus $49$ features).
SG-dyn used the full 1000-iteration budget in 6/7 selection folds on Glass,
7/10 on Diabetes and 10/10 on Abalone, and SG-const in 1/10 on Abalone
(DTAP in none); the voted weights of both SG variants are the best iterates
found within the budget. The SG results on Spambase, the only
dataset where they differ from DTAP, are given below Table~\ref{tab:acc};
Table~\ref{tab:perf-summary} includes all five methods.
Fold-wise stability can differ even when the final voted mask is identical.}

Table~\leavevmode\textcolor{red}{\ref{tab:wilcoxon} reports the signed-rank tests. None of the pairwise
comparisons rejects at $0.05$ after Holm adjustment; the SVM comparisons with
the two base models have adjusted $p$-values between $0.0625$ and $0.0938$.
The asymptotic Friedman tests reject equality globally, but this does not establish
a particular pairwise superiority. All DTAP--SG pairwise $p$-values equal $1$;
this reflects ties and limited evidence, not a statistical equivalence claim.
These small-sample comparisons concern eight heterogeneous datasets and do not
measure variation under repeated train/test splits.}

\begin{table}[htbp]
\centering\small\setlength{\tabcolsep}{4pt}
\caption{\textcolor{red}{Selected count $|A^\star|$ / Nogueira stability $\Psi$; $p$ is the original dimension. All five selection methods are shown.}}\label{tab:nfeat}
\begin{tabular}{lrrrrrr}
\toprule
Dataset & $p$ & N-Fscore & AFRSFCV & DTAP & \leavevmode\textcolor{red}{SG-const} & \leavevmode\textcolor{red}{SG-dyn} \\
\midrule
Glass & $9$ & \leavevmode\textcolor{red}{$7/0.816$} & \leavevmode\textcolor{red}{$7/0.609$} & \leavevmode\textcolor{red}{$7/0.633$} & \leavevmode\textcolor{red}{$7/0.633$} & \leavevmode\textcolor{red}{$7/0.633$} \\
Diabetes & $8$ & \leavevmode\textcolor{red}{$7/0.290$} & \leavevmode\textcolor{red}{$3/0.810$} & \leavevmode\textcolor{red}{$7/0.707$} & \leavevmode\textcolor{red}{$7/0.707$} & \leavevmode\textcolor{red}{$7/0.707$} \\
Segmentation & $18$ & \leavevmode\textcolor{red}{$8/0.885$} & $2/1.000$ & \leavevmode\textcolor{red}{$16/0.800$} & \leavevmode\textcolor{red}{$16/0.888$} & \leavevmode\textcolor{red}{$16/0.888$} \\
Abalone & $8$ & \leavevmode\textcolor{red}{$4/0.950$} & $2/1.000$ & $5/1.000$ & \leavevmode\textcolor{red}{$5/1.000$} & \leavevmode\textcolor{red}{$5/0.870$} \\
Satimage & $36$ & \leavevmode\textcolor{red}{$10/0.871$} & \leavevmode\textcolor{red}{$3/0.685$} & \leavevmode\textcolor{red}{$33/0.563$} & \leavevmode\textcolor{red}{$33/0.563$} & \leavevmode\textcolor{red}{$33/0.563$} \\
Yeast & $7$ & $6/1.000$ & $5/1.000$ & \leavevmode\textcolor{red}{$6/0.757$} & \leavevmode\textcolor{red}{$6/0.757$} & \leavevmode\textcolor{red}{$6/0.757$} \\
Spambase & $57$ & \leavevmode\textcolor{red}{$26/0.910$} & \leavevmode\textcolor{red}{$4/0.885$} & \leavevmode\textcolor{red}{$49/0.832$} & \leavevmode\textcolor{red}{$48/0.851$} & \leavevmode\textcolor{red}{$48/0.851$} \\
Pendigits & $16$ & \leavevmode\textcolor{red}{$13/0.854$} & \leavevmode\textcolor{red}{$12/0.063$} & $15/1.000$ & \leavevmode\textcolor{red}{$15/1.000$} & \leavevmode\textcolor{red}{$15/1.000$} \\
\bottomrule
\end{tabular}

\end{table}
\begin{table}[htbp]
\centering\small\setlength{\tabcolsep}{4pt}
\caption{\textcolor{red}{Per-dataset accuracy (\%) / macro one-vs-rest AUC on the single holdout. Equal voted masks reuse the same predictions; the SG exceptions are stated below.}}\label{tab:acc}
\begin{tabular}{llrrrr}
\toprule
Dataset & Classifier & Full & N-Fscore & AFRSFCV & DTAP \\
\midrule
Glass & SVM & \leavevmode\textcolor{red}{$76.74/0.9222$} & \leavevmode\textcolor{red}{$69.77/0.9073$} & \leavevmode\textcolor{red}{$72.09/0.9174$} & \leavevmode\textcolor{red}{$72.09/0.9174$} \\
 & $5$-NN & $72.09/0.9097$ & $74.42/0.8933$ & $69.77/0.9371$ & $69.77/0.9371$ \\
Diabetes & SVM & $74.68/0.8089$ & \leavevmode\textcolor{red}{$71.43/0.7957$} & $65.58/0.6617$ & $74.68/0.8098$ \\
 & $5$-NN & $71.43/0.7356$ & $70.78/0.7154$ & $67.53/0.7145$ & $73.38/0.7766$ \\
Segmentation & SVM & \leavevmode\textcolor{red}{$93.29/0.9936$} & \leavevmode\textcolor{red}{$80.52/0.9714$} & $78.79/0.9668$ & $93.29/0.9940$ \\
 & $5$-NN & $94.37/0.9880$ & $87.01/0.9801$ & $86.58/0.9646$ & $94.59/0.9897$ \\
Abalone & SVM & $61.60/0.8102$ & $61.36/0.8065$ & \leavevmode\textcolor{red}{$62.56/0.8114$} & $62.68/0.8123$ \\
 & $5$-NN & $59.81/0.7742$ & $61.00/0.7605$ & $58.73/0.7752$ & $60.29/0.7695$ \\
Satimage & SVM & $89.12/0.9796$ & \leavevmode\textcolor{red}{$87.10/0.9756$} & $80.42/0.9611$ & \leavevmode\textcolor{red}{$89.12/0.9797$} \\
 & $5$-NN & $91.14/0.9792$ & $87.80/0.9705$ & $80.96/0.9306$ & $91.22/0.9766$ \\
Yeast & SVM & \leavevmode\textcolor{red}{$58.25/0.8561$} & \leavevmode\textcolor{red}{$58.92/0.8241$} & \leavevmode\textcolor{red}{$57.91/0.8207$} & \leavevmode\textcolor{red}{$58.59/0.8426$} \\
 & $5$-NN & $51.52/0.7667$ & $51.52/0.7763$ & $51.52/0.7681$ & $51.52/0.7585$ \\
Spambase & SVM & \leavevmode\textcolor{red}{$92.07/0.9662$} & \leavevmode\textcolor{red}{$88.38/0.9392$} & $87.62/0.9155$ & \leavevmode\textcolor{red}{$91.97/0.9616$} \\
 & $5$-NN & $90.34/0.9443$ & $85.78/0.9214$ & $88.49/0.9206$ & $89.25/0.9437$ \\
Pendigits & SVM & $99.09/0.9995$ & $98.95/0.9996$ & $97.77/0.9996$ & $98.95/0.9995$ \\
 & $5$-NN & $98.95/0.9982$ & $98.64/0.9982$ & $97.73/0.9965$ & $99.05/0.9987$ \\
\bottomrule
\end{tabular}
\par\smallskip\begin{minipage}{0.97\textwidth}\footnotesize \leavevmode\textcolor{red}{Both SG variants equal DTAP exactly on the other seven datasets. On Spambase, both give $91.86/0.9628$ (SVM) and $89.14/0.9441$ ($5$-NN).} \end{minipage}
\end{table}
\begin{table}[htbp]
\centering\small\setlength{\tabcolsep}{4pt}
\caption{\textcolor{red}{Unweighted mean classification performance over the eight classification datasets. All five selection methods and the full-feature reference are included.}}\label{tab:perf-summary}
\begin{tabular}{lrrrr}
\toprule
Method & SVM accuracy & $5$-NN accuracy & SVM AUC & $5$-NN AUC \\
\midrule
Full & $80.61$ & $78.71$ & \leavevmode\textcolor{red}{$0.9170$} & $0.8870$ \\
N-Fscore & $77.05$ & $77.12$ & \leavevmode\textcolor{red}{$0.9024$} & $0.8770$ \\
AFRSFCV & $75.34$ & $75.16$ & \leavevmode\textcolor{red}{$0.8818$} & $0.8759$ \\
DTAP & $80.17$ & $78.63$ & \leavevmode\textcolor{red}{$0.9146$} & $0.8938$ \\
\leavevmode\textcolor{red}{SG-const} & \leavevmode\textcolor{red}{$80.16$} & \leavevmode\textcolor{red}{$78.62$} & \leavevmode\textcolor{red}{$0.9148$} & \leavevmode\textcolor{red}{$0.8939$} \\
\leavevmode\textcolor{red}{SG-dyn} & \leavevmode\textcolor{red}{$80.16$} & \leavevmode\textcolor{red}{$78.62$} & \leavevmode\textcolor{red}{$0.9148$} & \leavevmode\textcolor{red}{$0.8939$} \\
\bottomrule
\end{tabular}

\end{table}
\begin{table}[htbp]
\centering\footnotesize\setlength{\tabcolsep}{4pt}
\caption{\textcolor{red}{DTAP versus each baseline: raw / Holm-adjusted signed-rank $p$-values. Each metric has four comparisons over eight datasets. The last column gives the tie-corrected Friedman asymptotic $p$-value for all five methods.}}\label{tab:wilcoxon}
\begin{tabular}{lrrrrr}
\toprule
Metric & N-Fscore & AFRSFCV & SG-const & SG-dyn & Friedman \\
\midrule
Accuracy, SVM & \leavevmode\textcolor{red}{$0.0312/0.0938$} & \leavevmode\textcolor{red}{$0.0156/0.0625$} & $1.0000/1.0000$ & $1.0000/1.0000$ & \leavevmode\textcolor{red}{$0.0005$} \\
Accuracy, $5$-NN & \leavevmode\textcolor{red}{$0.3750/1.0000$} & \leavevmode\textcolor{red}{$0.0312/0.1250$} & $1.0000/1.0000$ & $1.0000/1.0000$ & \leavevmode\textcolor{red}{$0.0057$} \\
AUC, SVM & \leavevmode\textcolor{red}{$0.0156/0.0625$} & \leavevmode\textcolor{red}{$0.0312/0.0938$} & \leavevmode\textcolor{red}{$1.0000/1.0000$} & \leavevmode\textcolor{red}{$1.0000/1.0000$} & \leavevmode\textcolor{red}{$0.0055$} \\
AUC, $5$-NN & \leavevmode\textcolor{red}{$0.0781/0.3125$} & \leavevmode\textcolor{red}{$0.1562/0.4688$} & $1.0000/1.0000$ & $1.0000/1.0000$ & \leavevmode\textcolor{red}{$0.0435$} \\
\bottomrule
\end{tabular}

\end{table}

\subsection{\textcolor{red}{Constraint activity and a tuned subgradient comparison}}
\label{ssec:feasibility-audit}

\leavevmode\textcolor{red}{As noted in Section~\ref{ssec:reduction}, the calibrated thresholds need not yield
$\Omega\ne\varnothing$. For each full-data instance below, we keep the evaluated point with the largest slack
$\min_i(b_i-F_i(w))$ over all runs. Re-evaluated on matrices rebuilt from the data, each of
these points satisfies all three inequalities strictly; the smallest slack is
$4.33\times10^{-4}$. This is a numerical check, not an exact certificate of feasibility.}

\leavevmode\textcolor{red}{We count the iterates of a run, including the last (its evaluated points), at which
$F_3(w^k)-b_3>10^{-6}$. This measures violation along the trajectory, rather than whether the
constraint is binding at its limit. Table~\ref{tab:activation-tuning} shows that, from the
standard start, $F_3$ is violated at no evaluated point on seven of the eight classification
datasets, and at one of the $48$ evaluated points on Abalone. Since $F_3$ is convex
(Proposition~\ref{prop:fs-reg}), these feature-selection instances test a convex subclass and
do not demonstrate a need for the nonconvex extension. We also report the all-convex
comparator (AC), which assigns all three constraints to the convex track.}

\paragraph*{\textcolor{red}{Protocol and tuning cost.}}
\leavevmode\textcolor{red}{These are optimization-only tests on full-data feasibility instances.
They use $J\le10^{-6}$, at most $1\,000$ evaluated points, and identical
starting points across methods. SG-const and SG-dyn are defined in Section~\ref{sssec:baselines}. Besides SG-dyn with the
sampled modulus $\widehat L$, a tuned variant selects $L=c\widehat L$ from the manually specified grid
$c\in\{10^{-3},10^{-2.5},\ldots,10^1\}$ using three random starting
points (seeds $1101$--$1103$). Selection first maximizes the number of
successful runs, then minimizes total evaluations, then total best residual.
Evaluation uses the standard start and three different random starts
(seeds $2101$--$2103$). The feature starts are Dirichlet samples.
Only starting points are held out: tuning and evaluation use the same
full-data instance, and no classification score is used. Neither
$\widehat L$ nor the tuned value is a certified global H\"older modulus.}

\begin{table}[t]
\centering\small
\setlength{\tabcolsep}{4pt}
\caption{\textcolor{red}{Optimization evaluations from the uniform simplex start.
AC denotes the all-convex comparator; sampled and tuned refer to SG-dyn.
$\dagger$ means the tolerance was not reached within the budget.
The last column gives the DTAP third-constraint violation count divided
by its number of evaluated points.}}
\label{tab:activation-tuning}
\begin{tabular}{@{}lrrrrrl@{}}
\toprule
\leavevmode\textcolor{red}{Dataset} & \leavevmode\textcolor{red}{DTAP} & \leavevmode\textcolor{red}{AC} & \leavevmode\textcolor{red}{SG-const} & \leavevmode\textcolor{red}{Sampled} & \leavevmode\textcolor{red}{Tuned} & \leavevmode\textcolor{red}{$F_3$ activity}\\
\midrule
\leavevmode\textcolor{red}{Glass}        & \leavevmode\textcolor{red}{89} & \leavevmode\textcolor{red}{89} & \leavevmode\textcolor{red}{520} & \leavevmode\textcolor{red}{$1000^\dagger$} & \leavevmode\textcolor{red}{100} & \leavevmode\textcolor{red}{$0/89$}\\
\leavevmode\textcolor{red}{Diabetes}     & \leavevmode\textcolor{red}{14} & \leavevmode\textcolor{red}{14} & \leavevmode\textcolor{red}{350} & \leavevmode\textcolor{red}{$1000^\dagger$} & \leavevmode\textcolor{red}{51} & \leavevmode\textcolor{red}{$0/14$}\\
\leavevmode\textcolor{red}{Segmentation} & \leavevmode\textcolor{red}{3} & \leavevmode\textcolor{red}{3} & \leavevmode\textcolor{red}{124} & \leavevmode\textcolor{red}{170} & \leavevmode\textcolor{red}{3} & \leavevmode\textcolor{red}{$0/3$}\\
\leavevmode\textcolor{red}{Abalone}      & \leavevmode\textcolor{red}{48} & \leavevmode\textcolor{red}{48} & \leavevmode\textcolor{red}{638} & \leavevmode\textcolor{red}{$1000^\dagger$} & \leavevmode\textcolor{red}{3} & \leavevmode\textcolor{red}{$1/48$}\\
\leavevmode\textcolor{red}{Satimage}     & \leavevmode\textcolor{red}{2} & \leavevmode\textcolor{red}{2} & \leavevmode\textcolor{red}{60} & \leavevmode\textcolor{red}{102} & \leavevmode\textcolor{red}{2} & \leavevmode\textcolor{red}{$0/2$}\\
\leavevmode\textcolor{red}{Yeast}        & \leavevmode\textcolor{red}{4} & \leavevmode\textcolor{red}{4} & \leavevmode\textcolor{red}{324} & \leavevmode\textcolor{red}{305} & \leavevmode\textcolor{red}{3} & \leavevmode\textcolor{red}{$0/4$}\\
\leavevmode\textcolor{red}{Spambase}     & \leavevmode\textcolor{red}{5} & \leavevmode\textcolor{red}{5} & \leavevmode\textcolor{red}{93} & \leavevmode\textcolor{red}{295} & \leavevmode\textcolor{red}{56} & \leavevmode\textcolor{red}{$0/5$}\\
\leavevmode\textcolor{red}{Pendigits}    & \leavevmode\textcolor{red}{2} & \leavevmode\textcolor{red}{2} & \leavevmode\textcolor{red}{49} & \leavevmode\textcolor{red}{10} & \leavevmode\textcolor{red}{31} & \leavevmode\textcolor{red}{$0/2$}\\
\midrule
\leavevmode\textcolor{red}{Movement Libras} & \leavevmode\textcolor{red}{637} & \leavevmode\textcolor{red}{637} & \leavevmode\textcolor{red}{62} & \leavevmode\textcolor{red}{$1000^\dagger$} & \leavevmode\textcolor{red}{144} & \leavevmode\textcolor{red}{$0/637$}\\
\bottomrule
\end{tabular}
\end{table}

\leavevmode\textcolor{red}{Whenever $F_3(w^k)\le b_3$ along a run, the pseudoconvex track is empty,
$\mu_k=1$ and DTAP coincides step by step with the all-convex comparator;
on the eight instances with zero $F_3$ activity, the identical DTAP and AC
columns reflect this identity. On Abalone, $F_3$ is violated only at the uniform
start ($1/48$), where $J^{\ps}_0\approx1.56\times10^{-3}$ is negligible against
$J^{\cvx}_0\approx366.3$, so $1-\mu_0\approx4.27\times10^{-6}$; the values of $J$
along the two trajectories satisfy $|J_{\rm DTAP}-J_{\rm AC}|\le3.18\times10^{-5}\,J_{\rm AC}$ at
every evaluated point,
and both runs stop after $48$ evaluated points. The columns
are therefore not a comparison of the two assignments, whose success counts
differ only on the random starts ($26/36$ versus $36/36$}).

\leavevmode\textcolor{red}{Tuning raises SG-dyn's standard-start success count from $5/9$ to $9/9$.
It can outperform DTAP, for example on Abalone, and SG-const is faster
on Movement Libras. Across all four starts on all nine feature instances,
the success counts are $26/36$ for DTAP, $36/36$ for the all-convex
comparator, $33/36$ for SG-const, $15/36$ for sampled SG-dyn, and $28/36$
for tuned SG-dyn. A balanced DTAP variant also obtains $26/36$; it assigns
equal weights when both tracks are violated and full weight to the only violated track
otherwise. Thus these tests do not establish a uniform advantage of the proportional weight
or robustness to initialization.}

\leavevmode\textcolor{red}{The nine feature-instance searches required $179\,744$ optimization
evaluations and $10.80$ seconds of tuning solver time in total. These costs
are additional to Table~\ref{tab:activation-tuning}; DTAP was not tuned.
Initial modulus sampling, preprocessing and search overhead are excluded
from this solver-time total. The supplementary records give the cost of
each search separately.}

\paragraph*{\textcolor{red}{A genuinely nonconvex pseudoconvex constraint.}}
\leavevmode\textcolor{red}{To exercise the pseudoconvex track directly, consider $K=[-2,2]^d$ and}
\begin{equation}\label{eq:logball-test}\textcolor{red}{
\begin{aligned}
 F_1(x)&=\|x+e_1\|^2\le1,\qquad
 F_2(x)=\|x-e_2\|^2\le\tfrac94,\\
 F_3(x)&=\kappa\log(1+\|x-e_1\|^2)\le \kappa\log3.
\end{aligned}
}\end{equation}
\leavevmode\textcolor{red}{Here $d\in\{2,10,50\}$ and $\kappa\in\{0.3,1,3\}$ give nine configurations.
The third function is pseudoconvex: it is an increasing differentiable
transform, with positive derivative, of a convex quadratic. It is not
convex, since its radial Hessian eigenvalue is
$2\kappa(1-r^2)/(1+r^2)^2<0$ for $r=\|x-e_1\|>1$.
The exact witness $x=-e_1/10$ has positive slacks
$19/100$, $31/25$, and $\kappa\log(300/221)$.
The standard start is $1.5\mathbf1$; the three evaluation starts and three
tuning starts are drawn uniformly from $[0.8,1.8]^d$ with the seeds above
(NumPy \texttt{default\_rng}).
For modulus sampling we use $500$ uniform points in $K$ (seed $42$).
SG-const uses the same step $\nu_k=10^{-3}$ as on the feature instances. Tuned
SG-const uses a constant step $\nu$ selected from the grid
$\{10^{-3},10^{-2.5},\ldots,10^1\}$ with the tuning starts and selection rule of
tuned SG-dyn; it selects $\nu=10^{-0.5}$ in all nine configurations. The balanced
variant is the one described above.}

\begin{table}[t]
\centering\small
\setlength{\tabcolsep}{5pt}
\caption{\textcolor{red}{Nonconvex pseudoconvex tests, grouped by $\kappa$.
Each entry is successes out of $12$ runs, followed by the median number
of evaluated points in parentheses, including capped runs. Each group
contains all three dimensions and four starts. Tuned SG-const uses the step
selected as described in the text.}}
\label{tab:synthetic-scale}
\begin{tabular}{@{}lrrrrrr@{}}
\toprule
\leavevmode\textcolor{red}{$\kappa$} & \leavevmode\textcolor{red}{DTAP} & \leavevmode\textcolor{red}{Balanced} & \leavevmode\textcolor{red}{SG-const} & \leavevmode\textcolor{red}{Tuned SG-const} & \leavevmode\textcolor{red}{Sampled SG-dyn} & \leavevmode\textcolor{red}{Tuned SG-dyn}\\
\midrule
\leavevmode\textcolor{red}{$0.3$} & \leavevmode\textcolor{red}{$0\ (1000)$} & \leavevmode\textcolor{red}{$0\ (1000)$} & \leavevmode\textcolor{red}{$0\ (1000)$} & \leavevmode\textcolor{red}{$12\ (15)$} & \leavevmode\textcolor{red}{$8\ (725.5)$} & \leavevmode\textcolor{red}{$12\ (142.5)$}\\
\leavevmode\textcolor{red}{$1$}   & \leavevmode\textcolor{red}{$12\ (70)$} & \leavevmode\textcolor{red}{$12\ (395.5)$} & \leavevmode\textcolor{red}{$0\ (1000)$} & \leavevmode\textcolor{red}{$12\ (15)$} & \leavevmode\textcolor{red}{$12\ (301.5)$} & \leavevmode\textcolor{red}{$12\ (44)$}\\
\leavevmode\textcolor{red}{$3$}   & \leavevmode\textcolor{red}{$12\ (4)$} & \leavevmode\textcolor{red}{$12\ (4)$} & \leavevmode\textcolor{red}{$0\ (1000)$} & \leavevmode\textcolor{red}{$12\ (16)$} & \leavevmode\textcolor{red}{$12\ (167)$} & \leavevmode\textcolor{red}{$12\ (11)$}\\
\bottomrule
\end{tabular}
\end{table}

\leavevmode\textcolor{red}{DTAP encounters a third-constraint violation in every one of these $36$
runs. Its pooled activity is $12\,835/12\,877$ evaluated points.
Nevertheless, it reaches the tolerance in only $24/36$ runs, compared
with $32/36$ for sampled SG-dyn, $36/36$ for tuned SG-dyn and $36/36$ for tuned
SG-const (Table~\ref{tab:synthetic-scale}).
Tuned SG-const reaches the tolerance in every run with a median of $15$--$16$
evaluated points, fewer than any other method at $\kappa=0.3$ and $\kappa=1$; at
$\kappa=3$, DTAP and the balanced variant need a median of $4$. These counts
exclude the $147\,331$ evaluations of its tuning search.}

\leavevmode\textcolor{red}{With $\nu_k=10^{-3}$, SG-const moves at most $999\times10^{-3}=0.999$ from its start
within the budget, whereas every start lies at distance at least $1.11$ from
$\Omega$; its zero counts are therefore fixed by the step size and the budget.
For $\kappa\le1$ the convex track is satisfied from the second evaluated point on,
so DTAP then follows the single-track iteration of Corollary~\ref{cor:ps}; both tracks are
violated simultaneously at only $32$ of these $12\,877$ evaluated points.}

\leavevmode\textcolor{red}{All DTAP failures occur at $\kappa=0.3$. Changing $\kappa$ leaves the feasible set
unchanged, but $v_3$ and $\nabla F_3$ both scale with $\kappa$, so when
$\|g^{\ps}_k\|\le1$ the pseudoconvex step $(1-\mu_k)\gamma_k g^{\ps}_k$ scales with
$\gamma_k\kappa^2$. Table~\ref{tab:dtap-sweep} changes one DTAP parameter at a time
for $\kappa\in\{0.3,1\}$. With $\gamma_k=a(k+0.5)^{-0.55}$, $a=10$ and $a=100$ give
$12/12$ successes at $\kappa=0.3$; $\rho_{\cvx}=1$ gives $4/12$, all with $d=2$;
$\delta=10^{-1}$ reproduces the baseline iterates exactly (Remark~\ref{rem:delta-role}); the
faster-decaying steps $s=0.75$ and $s=1$ give no success at $\kappa=0.3$, and
$s=1$ also fails in all $12$ runs at $\kappa=1$. The failures at $\kappa=0.3$ are
therefore consistent with a diminishing step that is too small for this
constraint scale, rather than with a change in feasibility. These tests vary
one parameter at a time on one problem family and do not identify a setting
that is best in general.
On~\eqref{eq:logball-test}, $\delta=10^{-1}$ and $\delta=10^{-3}$ give identical iterates because
the pseudoconvex track has a single index, at every start either only $F_1$ is violated or
$v_1-v_2\ge4.79$, and the convex track is satisfied from the second evaluated point on; the
$\delta$ row of Table~\ref{tab:dtap-sweep} is therefore a consistency check, not evidence about
the role of $\delta$.
The synthetic tuning searches consumed $188\,295$ optimization evaluations and
$10.98$ seconds of solver time for SG-dyn, and $147\,331$ evaluations and $8.27$
seconds for SG-const, with the same exclusions as above, in the environment
stated below.}

\begin{table}[t]
\centering\small
\setlength{\tabcolsep}{5pt}
\caption{\textcolor{red}{DTAP on~\eqref{eq:logball-test} with one parameter changed at a time,
$\gamma_k=a(k+0.5)^{-s}$. The baseline $\rho_{\cvx}=1.98$, $\delta=10^{-3}$, $a=1$,
$s=0.55$ is the DTAP column of Table~\ref{tab:synthetic-scale}. Each entry is
successes out of $12$ runs (three dimensions, four starts), followed by the
median number of evaluated points in parentheses, including capped runs.}}
\label{tab:dtap-sweep}
\begin{tabular}{@{}lrr@{}}
\toprule
\leavevmode\textcolor{red}{Setting} & \leavevmode\textcolor{red}{$\kappa=0.3$} & \leavevmode\textcolor{red}{$\kappa=1$}\\
\midrule
\leavevmode\textcolor{red}{Baseline} & \leavevmode\textcolor{red}{$0\ (1000)$} & \leavevmode\textcolor{red}{$12\ (70)$}\\
\leavevmode\textcolor{red}{$\rho_{\cvx}=1$} & \leavevmode\textcolor{red}{$4\ (1000)$} & \leavevmode\textcolor{red}{$12\ (223.5)$}\\
\leavevmode\textcolor{red}{$\delta=10^{-1}$} & \leavevmode\textcolor{red}{$0\ (1000)$} & \leavevmode\textcolor{red}{$12\ (70)$}\\
\leavevmode\textcolor{red}{$a=10$} & \leavevmode\textcolor{red}{$12\ (77)$} & \leavevmode\textcolor{red}{$12\ (6)$}\\
\leavevmode\textcolor{red}{$a=100$} & \leavevmode\textcolor{red}{$12\ (6)$} & \leavevmode\textcolor{red}{$12\ (48)$}\\
\leavevmode\textcolor{red}{$s=0.75$} & \leavevmode\textcolor{red}{$0\ (1000)$} & \leavevmode\textcolor{red}{$12\ (374)$}\\
\leavevmode\textcolor{red}{$s=1$} & \leavevmode\textcolor{red}{$0\ (1000)$} & \leavevmode\textcolor{red}{$0\ (1000)$}\\
\bottomrule
\end{tabular}
\end{table}

\paragraph*{\textcolor{red}{Reproducibility.}}
\leavevmode\textcolor{red}{The feature-selection results were saved with Python~3.11.7 and NumPy~1.26.4
on Windows, with one BLAS/OpenMP thread requested. The feature oracles use algebraically
compressed quadratic forms; checked differences from the original
values and gradients are at most $1.82\times10^{-12}$ and
$1.10\times10^{-11}$, respectively. The solver times reported for the feature instances refer
to this compressed implementation. The accompanying package contains only the uncompressed
oracles; it reproduces the evaluation counts, while the reported solver times are those of
the original compressed implementation. Tables~\ref{tab:synthetic-scale}
and~\ref{tab:dtap-sweep} and the synthetic tuning costs were produced with
Python~3.13.4 and NumPy~2.1.3 on Windows with the default number of BLAS
threads; their counts are deterministic, and their solver times are single
measurements on that machine. With Python~3.13.4 and NumPy~2.1.3, the package also
reproduces exactly every numerical entry of Tables~\ref{tab:datasets}--\ref{tab:wilcoxon}
(Tables~\ref{tab:perf-summary} and~\ref{tab:wilcoxon} from unrounded scores) and of
Table~\ref{tab:activation-tuning}, and the counts quoted in Sections~\ref{ssec:results}
and~\ref{ssec:feasibility-audit}.
The supplementary scripts, raw trajectories, tuning records, witnesses,
source hashes and validation instructions accompany the results.}

\section{Conclusion}\label{sec:conclusion}
\leavevmode\textcolor{red}{DTAP solves normal-box nonlinear split feasibility problems with convex and pseudoconvex components by combining a residual-scaled convex step with a safeguarded pseudoconvex step. Under consistency and continuous differentiability, the whole sequence converges to a feasible point; no regularity modulus is an input, although the analysis uses bounds on a compact set containing the iterates. In the feature-selection application all constraints are convex, and test problems with a nonconvex pseudoconvex constraint exercise the pseudoconvex track. In the experiments, tuned subgradient methods can match or outperform DTAP, whose finite-budget performance depends on the starting point and the step scale. Rates for the full residual $J(x^k)$, beyond the convex-track bound of Remark~\ref{rem:rate}, and guarantees under error-bound conditions remain open.}

\section*{Declaration of competing interest}
The authors declare that they have no known competing financial interests or personal relationships that could have appeared to influence the work reported in this paper.
\section*{Funding}
The authors received no specific funding for this work.
\section*{Declaration of generative AI use}
\leavevmode\textcolor{red}{OpenAI Codex was used during revision for language editing.}
\section*{Data availability}
The datasets are available from the UCI Machine Learning Repository~\citep{kelly2023uci}. \leavevmode\textcolor{red}{The accompanying reproducibility package contains the scripts, recorded numerical outputs, and instructions for regenerating the reported tables.}

\section*{ORCID}
\noindent Nguyen Van The~\raisebox{-0.1ex}{\textit{https://orcid.org/0009-0004-6548-420X}}\\
Tran Ngoc Thang~\raisebox{-0.1ex}{\textit{https://orcid.org/0009-0000-9549-5736}}\\
Nguyen Hong Son~\raisebox{-0.1ex}{\textit{https://orcid.org/0009-0006-3106-8296}}


\begin{thebibliography}{99}
\newcommand{\bibdoi}[1]{doi: {\urlstyle{rm}\url{#1}}}

\bibitem[Bauschke and Combettes(2011)]{bauschke2011}
H.~H. Bauschke and P.~L. Combettes.
\newblock \emph{Convex Analysis and Monotone Operator Theory in {H}ilbert Spaces}.
\newblock CMS Books in Mathematics. Springer, New York, 2011.
\newblock \bibdoi{10.1007/978-1-4419-9467-7}.

\bibitem[Brown et~al.(2012)Brown, Pocock, Zhao, and Luj{\'a}n]{brown2012}
G.~Brown, A.~Pocock, M.-J.~Zhao, and M.~Luj{\'a}n.
\newblock Conditional likelihood maximization: a unifying framework for information theoretic feature selection.
\newblock \emph{Journal of Machine Learning Research}, 13:\penalty0 27--66, 2012.

\bibitem[Byrne(2002)]{byrne2002}
C.~Byrne.
\newblock Iterative oblique projection onto convex sets and the split feasibility problem.
\newblock \emph{Inverse Problems}, 18\penalty0 (2):\penalty0 441--453, 2002.
\newblock \bibdoi{10.1088/0266-5611/18/2/310}.

\bibitem[Censor and Elfving(1994)]{censor1994}
Y.~Censor and T.~Elfving.
\newblock A multiprojection algorithm using {B}regman projections in a product space.
\newblock \emph{Numerical Algorithms}, 8:\penalty0 221--239, 1994.
\newblock \bibdoi{10.1007/BF02142692}.

\bibitem[Combettes(2001)]{combettes2001quasifejer}
P.~L. Combettes.
\newblock Quasi-Fej{\'e}rian analysis of some optimization algorithms.
\newblock In D.~Butnariu, Y.~Censor, and S.~Reich, editors, \emph{Inherently Parallel Algorithms in Feasibility and Optimization and Their Applications}, volume~8 of \emph{Studies in Computational Mathematics}, pages 115--152. Elsevier, Amsterdam, 2001.
\newblock \bibdoi{10.1016/S1570-579X(01)80010-0}.

\bibitem[Condat(2016)]{condat2016simplex}
L.~Condat.
\newblock Fast projection onto the simplex and the $\ell_1$ ball.
\newblock \emph{Mathematical Programming}, 158\penalty0 (1):\penalty0 575--585, 2016.
\newblock \bibdoi{10.1007/s10107-015-0946-6}.

\bibitem[Hu et~al.(2022)Hu, Li, Yu, and Yip]{hu2022quasi}
Y.~Hu, G.~Li, C.~K.~W. Yu, and T.~L. Yip.
\newblock Quasi-convex feasibility problems: Subgradient methods and convergence rates.
\newblock \emph{European Journal of Operational Research}, 298\penalty0 (1):\penalty0 45--58, 2022.
\newblock \bibdoi{10.1016/j.ejor.2021.09.029}.

\bibitem[Kelly et~al.(2023)Kelly, Longjohn, and Nottingham]{kelly2023uci}
M.~Kelly, R.~Longjohn, and K.~Nottingham.
\newblock The {UCI} machine learning repository.
\newblock \url{https://archive.ics.uci.edu}, 2023.
\newblock Accessed 2026.

\bibitem[Li et~al.(2013)Li, Han, and Zhang]{li2013nsfp}
Z.~Li, D.~Han, and W.~Zhang.
\newblock A self-adaptive projection-type method for nonlinear multiple-sets split feasibility problem.
\newblock \emph{Inverse Problems in Science and Engineering}, 21\penalty0 (1):\penalty0 155--170, 2013.
\newblock \bibdoi{10.1080/17415977.2012.677445}.

\bibitem[L{\'o}pez et~al.(2012)L{\'o}pez, Mart{\'\i}n-M{\'a}rquez, Wang, and Xu]{lopez2012}
G.~L{\'o}pez, V.~Mart{\'\i}n-M{\'a}rquez, F.~Wang, and H.-K. Xu.
\newblock Solving the split feasibility problem without prior knowledge of matrix norms.
\newblock \emph{Inverse Problems}, 28\penalty0 (8):\penalty0 085004, 2012.
\newblock \bibdoi{10.1088/0266-5611/28/8/085004}.

\bibitem[Mangasarian(1965)]{mangasarian1965}
O.~L. Mangasarian.
\newblock Pseudo-convex functions.
\newblock \emph{SIAM Journal on Control}, 3\penalty0 (2):\penalty0 281--290, 1965.
\newblock \bibdoi{10.1137/0303020}.

\bibitem[Mangasarian(1994)]{mangasarian1994}
O.~L. Mangasarian.
\newblock \emph{Nonlinear Programming}.
\newblock Classics in Applied Mathematics. SIAM, Philadelphia, 1994.
\newblock \bibdoi{10.1137/1.9781611971255}.

\bibitem[Nogueira et~al.(2018)Nogueira, Sechidis, and Brown]{nogueira2018}
S.~Nogueira, K.~Sechidis, and G.~Brown.
\newblock On the stability of feature selection algorithms.
\newblock \emph{Journal of Machine Learning Research}, 18\penalty0 (174):\penalty0 1--54, 2018.

\bibitem[Polyak(1969)]{polyak1969}
B.~T. Polyak.
\newblock Minimization of unsmooth functionals.
\newblock \emph{USSR Computational Mathematics and Mathematical Physics}, 9\penalty0 (3):\penalty0 14--29, 1969.
\newblock \bibdoi{10.1016/0041-5553(69)90061-5}.

\bibitem[Wang et~al.(2021)Wang, Li, and Wang]{wang2021fractional}
Y.~Wang, X.~Li, and J.~Wang.
\newblock A neurodynamic optimization approach to supervised feature selection via fractional programming.
\newblock \emph{Neural Networks}, 136:\penalty0 194--206, 2021.
\newblock \bibdoi{10.1016/j.neunet.2021.01.004}.

\bibitem[Wang et~al.(2026)Wang, Zhang, An, and Deng]{wang2026afrsfcv}
C.~Wang, Y.~Zhang, S.~An, and T.~Deng.
\newblock Adaptive feature selection based on fuzzy rough set fusion model with class variance.
\newblock \emph{Pattern Recognition}, 170:\penalty0 112014, 2026.
\newblock \bibdoi{10.1016/j.patcog.2025.112014}.

\bibitem[Yang(2004)]{yang2004}
Q.~Yang.
\newblock The relaxed {CQ} algorithm solving the split feasibility problem.
\newblock \emph{Inverse Problems}, 20\penalty0 (4):\penalty0 1261--1266, 2004.
\newblock \bibdoi{10.1088/0266-5611/20/4/014}.

\bibitem[Yang et~al.(2014)Yang, Yang, and Zhang]{yang2014alternating}
Y.~Yang, Q.~Yang, and S.~Zhang.
\newblock Modified alternating direction methods for the modified multiple-sets split feasibility problems.
\newblock \emph{Journal of Optimization Theory and Applications}, 163\penalty0 (1):\penalty0 130--147, 2014.
\newblock \bibdoi{10.1007/s10957-013-0502-6}.

\bibitem[Yen et~al.(2019)Yen, Huyen, and Muu]{yen2019}
L.~H. Yen, N.~T.~T. Huyen, and L.~D. Muu.
\newblock A subgradient algorithm for a class of nonlinear split feasibility problems: application to jointly constrained Nash equilibrium models.
\newblock \emph{Journal of Global Optimization}, 73:\penalty0 849--868, 2019.
\newblock \bibdoi{10.1007/s10898-018-00735-0}.

\bibitem[Zhang et~al.(2009)Zhang, Han, and Li]{zhang2009adaptive}
W.~Zhang, D.~Han, and Z.~Li.
\newblock A self-adaptive projection method for solving the multiple-sets split feasibility problem.
\newblock \emph{Inverse Problems}, 25\penalty0 (11):\penalty0 115001, 2009.
\newblock \bibdoi{10.1088/0266-5611/25/11/115001}.

\bibitem[Zhang et~al.(2012)Zhang, Han, and Yuan]{zhang2012simultaneous}
W.~Zhang, D.~Han, and X.~Yuan.
\newblock An efficient simultaneous method for the constrained multiple-sets split feasibility problem.
\newblock \emph{Computational Optimization and Applications}, 52\penalty0 (3):\penalty0 825--843, 2012.
\newblock \bibdoi{10.1007/s10589-011-9429-8}.

\bibitem[Zhao and Yang(2011)]{zhao2011adaptive}
J.~Zhao and Q.~Yang.
\newblock Self-adaptive projection methods for the multiple-sets split feasibility problem.
\newblock \emph{Inverse Problems}, 27\penalty0 (3):\penalty0 035009, 2011.
\newblock \bibdoi{10.1088/0266-5611/27/3/035009}.

\end{thebibliography}

\end{document}
